% =====================================================================
%  U15 —— 第 10 章 10.3 无分歧情形
%  原文 PDF p.326–p346（印刷 p.312–332）
%  本文件不含 \chapter，以 \section 开头（接在第 10 章内）
% =====================================================================

\section{无分歧情形}

\begin{keypoints}
    \item \textbf{核心思想：} 无分歧情形是几何 Langlands 对应的最简情形：唯一的无分歧局部系统是平凡局部系统 $\sigma_0$，与它对应的 opers 构成不交并 $\bigsqcup_{\lambda\in P^+}\operatorname{Op}^\lambda_{{}^LG}$。
    \item \textbf{数学工具：} 球面 Hecke 代数 $H(G(F),K_0)$ 与 Satake 同构、Harish--Chandra 模 $(\hat{\mathfrak{g}}_{\kappa_c},G[[t]])$-mod、中心 $\mathfrak{z}=\operatorname{Fun}\operatorname{Op}_{{}^LG}(D)$ 上的模。
    \item \textbf{关键结论：} 范畴 $\hat{\mathfrak{g}}_{\kappa_c}\text{-mod}^{G[[t]]}_\chi$ 等价于向量空间范畴（定理 10.3.3）；更一般地 $\hat{\mathfrak{g}}_{\kappa_c}\text{-mod}^{G[[t]]}_{reg}$ 等价于 $\mathfrak{z}$-mod（定理 10.3.4）。
    \item \textbf{关键结论：} 球面 Hecke 代数的作用在几何情形由 Hecke 范畴给出；Hecke 本征向量条件提升为 Hecke 本征模条件 (10.3-26)。
    \item \textbf{记号提醒：} $P^+$ 为 $G$ 的整优权（等价地 ${}^LG$ 的整优余权）；$\mathfrak{z}=\mathfrak{z}(\hat{\mathfrak{g}})=\operatorname{Fun}\operatorname{Op}_{{}^LG}(D)$；$\operatorname{Op}^0_{{}^LG}=\operatorname{Op}_{{}^LG}(D)$。
\end{keypoints}

我们首先处理 $G((t))$ 的“极大紧子群” $K_0=G[[t]]$ 的情形，考察那些含有非平凡 $K_0$-等变对象的范畴 $\hat{\mathfrak{g}}_{\kappa_c}\text{-mod}_\chi$。

\subsection{$G(F)$ 的无分歧表示}

这些范畴是群 $G(F)$ 的光滑表示的类比，其中 $F$ 是局部非阿基米德域（例如 $\mathbb{F}_q((t))$），而这些光滑表示含有非零的 $K_0$-不变向量。这样的表示称为\textbf{无分歧的}。$G(F)$ 的不可约无分歧表示的分类是局部 Langlands 对应（见 1.1.4 与 1.1.5 节）中最简单的情形。具体地说，下列对象的等价类集合之间存在双射：
\begin{equation}
    \boxed{\begin{array}{c}\text{无分歧容许}\\[-2pt]\text{同态 } W'_F\to{}^LG\end{array}}
    \quad\Longleftrightarrow\quad
    \boxed{\begin{array}{c}\text{$G(F)$ 的不可约}\\[-2pt]\text{无分歧表示}\end{array}}
    \tag{10.3-1}
\end{equation}
其中 $W'_F$ 是 1.1.3 节引入的 Weil--Deligne 群。

% [Tutor Note: (10.3-1)–(10.3-4) 的逻辑链须记清：Satake 同构 (10.3-3) 把球面 Hecke 代数与 $\operatorname{Rep}{}^LG$ 等同；$\operatorname{Rep}{}^LG$ 交换，故其不可约模都是一维，其特征标由 ${}^LG$ 的半单共轭类参数化，得 (10.3-4)；再由“无分歧表示 ↔ 不可约 $H$-模”得 (10.3-1)、(10.3-2)。这里 ${}^LG$ 是复化群，其半单元素的共轭类构成复代数簇。]

按定义，无分歧同态 $W'_F\to{}^LG$ 指通过商
\[
    W'_F\to W_F\to\mathbb{Z}
\]
分解的同态（这些群与同态的定义见 1.1.3 节）。其容许性指它在 ${}^LG$ 中的像由半单元素组成。因此 (10.3-1) 左端的集合就是 ${}^LG$ 半单元素的共轭类集合。于是上面的双射可以重新解释为：
\begin{equation}
    \boxed{\begin{array}{c}\text{${}^LG$ 中的半单}\\[-2pt]\text{共轭类}\end{array}}
    \quad\Longleftrightarrow\quad
    \boxed{\begin{array}{c}\text{$G(F)$ 的不可约}\\[-2pt]\text{无分歧表示}\end{array}}
    \tag{10.3-2}
\end{equation}

为构造这个双射，我们来看\textbf{球面 Hecke 代数}\transnote{这里的“球面”指双 $K_0$-不变性：$H(G(F),K_0)$ 以双陪集 $K_0\backslash G(F)/K_0$ 的特征函数为基，乘法是卷积 $(f_1*f_2)(g)=\int_{G(F)}f_1(h)f_2(h^{-1}g)\,dh$。Satake 同构 (10.3-3) 说这个卷积代数是交换的，并同构于 $\operatorname{Rep}{}^LG$；正因交换，它的不可约模都是一维，才有无分歧表示与半单共轭类的双射。后文用 Iwahori 子群 $I$ 代替 $K_0$ 得到的 Iwahori Hecke 代数不再交换，这是“球面”一词用于区分两种 Hecke 结构的原因。[来源：原书 p.316；I. Satake, Theory of spherical functions on reductive algebraic groups over $p$-adic fields, Publ. Math. IHÉS 18 (1963), 5--69.]}
\[
    H(G(F),K_0).
\]
由 Satake 同构 [Satake 1963]（在 Langlands [Langlands 1970] 的解释下），此代数是交换的，且同构于 Langlands 对偶群 ${}^LG$ 的表示环：
\begin{equation}
    H(G(F),K_0)\simeq\operatorname{Rep}{}^LG.
    \tag{10.3-3}
\end{equation}
我们回忆，$\operatorname{Rep}{}^LG$ 由有限线性组合 $\sum_i a_i[V_i]$ 组成，其中 $V_i$ 是 ${}^LG$ 的有限维表示（不失一般性可设它们不可约），$a_i\in\mathbb{C}$；乘法为
\[
    [V]\cdot[W]=[V\otimes W].
\]
由于 $\operatorname{Rep}{}^LG$ 交换，它的不可约模都是一维的。它们对应特征标 $\operatorname{Rep}{}^LG\to\mathbb{C}$。有双射
\begin{equation}
    \boxed{\begin{array}{c}\text{${}^LG$ 中的半单}\\[-2pt]\text{共轭类}\end{array}}
    \quad\Longleftrightarrow\quad
    \boxed{\begin{array}{c}\operatorname{Rep}{}^LG\\[-2pt]\text{的特征标}\end{array}}
    \tag{10.3-4}
\end{equation}
其中对应于共轭类 $\gamma$ 的特征标 $\phi_\gamma$ 由公式
\[
    \phi_\gamma:[V]\mapsto\operatorname{Tr}(\gamma,V)
\]
给出。\footnote{通常会把此公式的右端（对不可约表示 $V$）乘以一个依赖于 $q$ 与 $V$ 的最高权的标量，但这对我们的讨论无关紧要。}

现在，若 $(R,\pi)$ 是 $G(F)$ 的表示，则 $V$ 中 $K_0$-不变向量构成的空间 $R^{\pi(K_0)}$ 是 $H(G(F),K_0)$ 上的模。容易证明，这给出了 $G(F)$ 的不可约无分歧表示的等价类与不可约 $H(G(F),K_0)$-模之间的一一对应。把它与双射 (10.3-4) 及同构 (10.3-3) 结合，就得到所求的双射 (10.3-1) 与 (10.3-2)。

% [Tutor Note: 原文此处写 “$K_0$-invariant vectors in $V$”，但本段以 $(R,\pi)$ 记表示、下文一律用 $R^{\pi(K_0)}$，故 $V$ 应是 $R$ 之误。译文保留原样，存疑待考。]

特别地，我们发现：由于 Hecke 代数 $H(G(F),K_0)$ 交换，不可约表示的 $K_0$-不变向量空间 $R^{\pi(K_0)}$（它是不可约的 $H(G(F),K_0)$-模）要么是 $0$，要么是一维的。若是一维的，则 $H(G(F),K_0)$ 按某个 $\gamma$ 的特征标 $\phi_\gamma$ 作用于它：
\begin{equation}
    H_V\star v=\phi_\gamma([V])v=\operatorname{Tr}(\gamma,V)v,
    \qquad v\in R^{\pi(K_0)},\ [V]\in\operatorname{Rep}{}^LG,
    \tag{10.3-5}
\end{equation}
其中 $H_V$ 是同构 (10.3-3) 下对应于 $[V]$ 的 $H(G(F),K_0)$ 中元素（卷积作用的定义见公式 (10.2-2)）。于是，任何 $v\in R^{\pi(K_0)}$ 都是 \textbf{Hecke 本征向量}。

% [Tutor Note: 式 (10.3-5) 的归一化依赖脚注 3 中被略去的标量：若与 1.1 节的局部 Langlands 对应比较，须把 $\operatorname{Tr}(\gamma,V)$ 乘以依赖于 $q$ 与 $V$ 最高权的因子。本小节的论证不依赖此归一化，但引用 (10.3-5) 作数值计算时不可忽略。]

现在我们讨论这些陈述的范畴类比。

\subsection{$\hat{\mathfrak{g}}_{\kappa_c}$-模的无分歧范畴}

在范畴情形，$G(F)$ 的不可约表示 $(R,\pi)$ 所扮演的角色，由某个 $\chi\in\operatorname{Op}_{{}^LG}(D^\times)$ 对应的范畴 $\hat{\mathfrak{g}}_{\kappa_c}\text{-mod}_\chi$ 担任。无分歧表示的类比，是含有非零 $(\hat{\mathfrak{g}}_{\kappa_c},G[[t]])$ Harish--Chandra 模的范畴 $\hat{\mathfrak{g}}_{\kappa_c}\text{-mod}_\chi$。这引出如下问题：对哪些 $\chi\in\operatorname{Op}_{{}^LG}(D^\times)$，范畴 $\hat{\mathfrak{g}}_{\kappa_c}\text{-mod}_\chi$ 含有非零的 $(\hat{\mathfrak{g}}_{\kappa_c},G[[t]])$ Harish--Chandra 模？

上一节我们看到，$(R,\pi)$ 无分歧当且仅当它对应于无分歧 Langlands 参数，即通过 $W'_F\to\mathbb{Z}$ 分解的同态 $W'_F\to{}^LG$。回忆在几何情形中，Langlands 参数是 $D^\times$ 上的 ${}^LG$-局部系统。无分歧同态 $W'_F\to{}^LG$ 的类比，是那些可延拓到圆盘 $D$ 的 $D^\times$ 上的局部系统；换言之，是在原点 $0\in D$ 处没有奇点的局部系统。注意 $D$ 上的局部系统在同构意义下唯一。事实上，设给定 $D$ 上 ${}^LG$-丛 $\mathcal{F}$ 上的正则联络。把 $\mathcal{F}$ 在 $0\in D$ 处的纤维 $\mathcal{F}_0$ 平凡化。则由于 $D$ 可缩，该联络把 $\mathcal{F}$ 与 $D$ 上的平凡丛等同。在此等同下，联络本身变为平凡联络，即由算子 $\nabla=\partial_t$ 表示。

因此，所有正则 ${}^LG$-局部系统（即那些可延拓到 $D$ 的局部系统）都对应于集合 $\operatorname{Loc}_{{}^LG}(D^\times)$ 的同一个点，即平凡局部系统 $\sigma_0$ 的等价类。\footnote{但要注意，$D$ 上的平凡 ${}^LG$-局部系有非平凡的自动同构群，即群 ${}^LG$ 本身（它可以实现为 $0\in D$ 处纤维的自动同构群）。因此若把 $\operatorname{Loc}_{{}^LG}(D^\times)$ 看作栈而非集合，则平凡局部系对应子栈 $\mathrm{pt}/{}^LG$。}从把 $\operatorname{Loc}_{{}^LG}(D^\times)$ 实现为商 (1.2-7) 的观点看，这只意味着：存在唯一的 ${}^LG((t))$ 规范等价类，它包含所有形如 $\partial_t+A(t)$（$A(t)\in{}^L\mathfrak{g}[[t]]$）的正则联络。

正则联络的规范等价类，就是几何情形中唯一可视为无分歧的局部 Langlands 参数。因此，类比于 $G(F)$ 的无分歧 Langlands 对应，我们预期：范畴 $\hat{\mathfrak{g}}_{\kappa_c}\text{-mod}_\chi$ 含有非零 $(\hat{\mathfrak{g}}_{\kappa_c},G[[t]])$ Harish--Chandra 模，当且仅当 ${}^LG$-oper $\chi\in\operatorname{Op}_{{}^LG}(D^\times)$ 在 ${}^LG((t))$ 的规范变换下等价于平凡联络；换言之，$\chi$ 属于 $\sigma_0$ 上的纤维 $\alpha^{-1}(\sigma_0)$。

% [Tutor Note: 这里 $\alpha:\operatorname{Op}_{{}^LG}(D^\times)\to\operatorname{Loc}_{{}^LG}(D^\times)$ 是 opers 到其底层局部系统的映射（定义见 9.2 节）。关键点：经典情形中无分歧参数只有一个共轭类，而 $\alpha^{-1}(\sigma_0)$ 并非单点，而是整整一族 opers，即引理 10.3.1 中的 $\bigsqcup_{\lambda\in P^+}\operatorname{Op}^\lambda_{{}^LG}$。这个“一个局部系统、多个 opers”的现象正是 10.3.7 节群胚作用的来源。]

这个纤维长什么样？设 $P^+$ 是 $G$ 的整优权集合（等价地，${}^LG$ 的整优余权集合）。在 9.2.3 节中，对每个 $\lambda\in P^+$ 定义了空间 $\operatorname{Op}^\lambda_{{}^LG}$，它由形如
\begin{equation}
    \nabla=\partial_t+\sum_{i=1}^{\ell}
    t^{\langle\check\alpha_i,\lambda\rangle}\psi_i(t)f_i+\mathbf{v}(t)
    \tag{10.3-6}
\end{equation}
的算子的 ${}^LB[[t]]$-等价类组成，其中 $\psi_i(t)\in\mathbb{C}[[t]]$，$\psi_i(0)\neq0$，$\mathbf{v}(t)\in{}^L\mathfrak{b}[[t]]$。

\begin{lemma}
    设 oper $\chi\in\operatorname{Op}_{{}^LG}(D^\times)$ 底层的局部系统平凡。则 $\chi$ 属于子集
    \[
        \operatorname{Op}^\lambda_{{}^LG}\subset\operatorname{Op}_{{}^LG}(D^\times),
        \qquad\lambda\in P^+,
    \]
    的不交并。
\end{lemma}

\begin{proof}
    由定义显然，$\operatorname{Op}^\lambda_{{}^LG}$ 中任何 oper 在圆盘 $D$ 上正则，因此在 ${}^LG((t))$ 的规范变换下等价于平凡联络。反过来，设 oper $\chi=(\mathcal{F},\nabla,\mathcal{F}_{{}^LB})$ 的底层 ${}^LG$-局部系统平凡。则 $\nabla$ 在 ${}^LG((t))$ 的规范变换下等价于某个正则联络，即形如 $\partial_t+A(t)$，其中 $A(t)\in{}^L\mathfrak{g}[[t]]$。我们有分解 ${}^LG((t))={}^LG[[t]]{}^LB((t))$。${}^LG[[t]]$ 的规范作用显然保持正则联络构成的空间。因此，若 oper 联络 $\nabla$ 在 ${}^LG((t))$ 的规范变换下等价于某个正则联络，则其 ${}^LB((t))$ 规范类中已经必须含有一个正则联络。oper 条件于是蕴含这个规范类含有形如 (10.3-6) 的联络算子，其中 $\lambda$ 是 ${}^LG$ 的某个整优权。因此 $\chi\in\operatorname{Op}^\lambda_{{}^LG}$。
\end{proof}

于是我们看到，对应于（唯一的）无分歧 Langlands 参数的 opers 集合，就是不交并 $\bigsqcup_{\lambda\in P^+}\operatorname{Op}^\lambda_{{}^LG}$。我们称这样的 opers 为\textbf{无分歧的}\transnote{“无分歧 oper”指其底层局部系统平凡、因而不含单值性；文献中也称这类对象为“无单值性的 oper”（opers without monodromy）。注意 $\alpha^{-1}(\sigma_0)$ 并非单点：$\lambda=0$ 给出在 $0$ 处正则的 opers，$\lambda\neq0$ 的 $\operatorname{Op}^\lambda$ 则在 $0$ 处有由 $\lambda$ 决定的极点阶数，但它们的底层局部系统仍平凡。$\lambda$ 记录的是 oper 的 ${}^LB$-约化（相对位置），而非局部系统的同构类；这正是“一个局部系统、多个 opers”的来源。[来源：原书 p.318--319；E. Frenkel and D. Gaitsgory, Weyl modules and opers without monodromy, arXiv:math.QA/0706.3725.]}。下面的结果证实了我们的预期：对应于无分歧 oper $\chi$ 的范畴 $\hat{\mathfrak{g}}_{\kappa_c}\text{-mod}_\chi$ 也“无分歧”，即含有非零 $G[[t]]$-等变对象，当且仅当 $\chi$ 无分歧。

\begin{lemma}
    范畴 $\hat{\mathfrak{g}}_{\kappa_c}\text{-mod}_\chi$ 含有非零 $(\hat{\mathfrak{g}}_{\kappa_c},G[[t]])$ Harish--Chandra 模，当且仅当
    \begin{equation}
        \chi\in\bigcup_{\lambda\in P^+}\operatorname{Op}^\lambda_{{}^LG}.
        \tag{10.3-7}
    \end{equation}
\end{lemma}

\begin{proof}
    记 $G[[t]]=G\ltimes K_1$，其中 $K_1$ 是 $G[[t]]$ 的第一同余子群，为 pro-幂幺的 pro-代数群。$\hat{\mathfrak{g}}_{\kappa_c}$-模 $M$ 是 $G[[t]]$-等变的，当且仅当 $K_1$ 的李代数 $\mathfrak{g}\otimes t\mathbb{C}[[t]]$ 在 $M$ 上的作用是局部幂零的、可以指数化提升为 $K_1$ 的作用，并且 $M$ 在李子代数 $\mathfrak{g}\otimes1$ 的作用下分解为有限维表示的直和。

    设范畴 $\hat{\mathfrak{g}}_{\kappa_c}\text{-mod}_\chi$ 含有非零的 $(\hat{\mathfrak{g}}_{\kappa_c},G[[t]])$ Harish--Chandra 模 $M$。由于 $M$ 光滑，任何向量 $v\in M$ 都被某个 $\mathfrak{g}\otimes t^N\mathbb{C}[[t]]$（$N>0$）零化。考虑 $M$ 的 $\mathfrak{g}\otimes t\mathbb{C}[[t]]$-子模 $M_v=U(\mathfrak{g}\otimes t\mathbb{C}[[t]])v$。由构造，$\mathfrak{g}\otimes t\mathbb{C}[[t]]$ 在 $M_v$ 上的作用通过商 $(\mathfrak{g}\otimes t\mathbb{C}[[t]])/(\mathfrak{g}\otimes t^N\mathbb{C}[[t]])$ 分解。由于 $M$ 是 $G[[t]]$-等变的，可知这个李代数的作用可以指数化为相应的代数群 $K_1/K_N$ 的代数作用，其中 $K_N$ 是 $G[[t]]$ 的第 $N$ 个同余子群，它是 $K_1$ 的正规子群。群 $K_1/K_N$ 是有限维幂幺代数群，于是得到 $K_1/K_N$ 在 $M_v$ 上的代数表示。任何这样的表示都含有不变向量。\footnote{这可由加法群 $\mathbb{G}_a$ 的相应命题归纳得到。}于是 $M_v$（从而 $M$）含有被 $\mathfrak{g}\otimes t\mathbb{C}[[t]]$ 零化的向量 $v'$。

    由 $v'$ 生成的 $M$ 的 $\mathfrak{g}[[t]]$-子模（等价地，$\mathfrak{g}$-子模）是 $\mathfrak{g}$ 的不可约有限维表示的直和。取其中一个不可约表示，记作 $M'$。由构造，$M'$ 被 $\mathfrak{g}\otimes t\mathbb{C}[[t]]$ 零化，且作为 $\mathfrak{g}$-模同构于 $V_\lambda$，即最高权为 $\lambda$ 的 $\mathfrak{g}$ 的不可约表示。因此存在非零同态 $\mathbb{V}_\lambda\to M$，其中 $\mathbb{V}_\lambda$ 是 9.6 节引入的 Weyl 模，它把生成子空间 $V_\lambda\subset\mathbb{V}_\lambda$ 送到 $M'$。

    按假设，$M$ 是范畴 $\hat{\mathfrak{g}}_{\kappa_c}\text{-mod}_\chi$ 的对象，中心 $Z(\hat{\mathfrak{g}})$ 按对应于 $\chi\in\operatorname{Op}_{{}^LG}(D^\times)$ 的中心特征标作用。非零同态 $\mathbb{V}_\lambda\to M$ 的存在，蕴含 $\chi$ 属于 $Z(\hat{\mathfrak{g}})$ 在 $\operatorname{End}_{\hat{\mathfrak{g}}_{\kappa_c}}\mathbb{V}_\lambda$ 中的像的谱。由定理 9.6.1，后者的谱等于 $\operatorname{Op}^\lambda_{{}^LG}$。因此得到 $\chi\in\operatorname{Op}^\lambda_{{}^LG}$。

    反之，若 $\chi\in\operatorname{Op}^\lambda_{{}^LG}$，则由定理 9.6.1，$\mathbb{V}_\lambda$ 对相应于 $\chi$ 的中心特征标的商非零。这个商显然是范畴 $\hat{\mathfrak{g}}_{\kappa_c}\text{-mod}_\chi$ 的 $G[[t]]$-等变对象。证毕。
\end{proof}

% [Tutor Note: 证明中 $G[[t]]=G\ltimes K_1$ 的 $K_1$ 是 pro-幂幺 pro-代数群；$K_0$-等变性要求 $\mathfrak{g}\otimes t\mathbb{C}[[t]]$ 局部幂零（来自光滑性）且 $\mathfrak{g}\otimes 1$ 的作用可分解（来自代数性）。这两条缺一不可，是 Harish--Chandra 模定义中“$(\hat{\mathfrak{g}}_{\kappa_c},G[[t]])$”配对的实际含义。]

下一个问题是描述 $\chi\in\operatorname{Op}^\lambda_{{}^LG}$ 时 $(\hat{\mathfrak{g}}_{\kappa_c},G[[t]])$ Harish--Chandra 模构成的范畴 $\hat{\mathfrak{g}}_{\kappa_c}\text{-mod}^{G[[t]]}_\chi$。

\subsection{$G[[t]]$-等变模的范畴}

回忆 10.3.1 节：$G(F)$ 的无分歧不可约表示中，$K_0$-不变向量构成的空间总是一维的。按我们的方案，范畴 $\hat{\mathfrak{g}}_{\kappa_c}\text{-mod}^{G[[t]]}_\chi$ 应被看作这个空间的范畴类比。因此我们预期它是最简单的阿贝尔范畴：$\mathbb{C}$-向量空间范畴。回忆这里假设 $\chi$ 属于空间 $\operatorname{Op}^\lambda_{{}^LG}$ 的并，其中 $\lambda\in P^+$；否则范畴 $\hat{\mathfrak{g}}_{\kappa_c}\text{-mod}^{G[[t]]}_\chi$ 将是平凡范畴（即零对象是唯一对象）。

本小节我们将按 [Frenkel and Gaitsgory 2004]（另见 [Beilinson and Drinfeld 1997]）证明：取 $\lambda=0$ 时（此时 $\operatorname{Op}^0_{{}^LG}=\operatorname{Op}_{{}^LG}(D)$，故 $\chi\in\operatorname{Op}_{{}^LG}(D)\subset\operatorname{Op}_{{}^LG}(D^\times)$）预期确实成立。我们还将猜想：对一切 $\lambda\in P^+$ 与 $\chi\in\operatorname{Op}^\lambda_{{}^LG}$，这也成立。

回忆真空模 $\mathbb{V}_0=V_{\kappa_c}(\mathfrak{g})$。由同构 (9.5-1)，我们有
\begin{equation}
    \operatorname{End}_{\hat{\mathfrak{g}}_{\kappa_c}}\mathbb{V}_0\simeq\operatorname{Fun}\operatorname{Op}_{{}^LG}(D).
    \tag{10.3-8}
\end{equation}
设 $\chi\in\operatorname{Op}_{{}^LG}(D)\subset\operatorname{Op}_{{}^LG}(D^\times)$。则 $\chi$ 定义了代数 $\operatorname{End}_{\hat{\mathfrak{g}}_{\kappa_c}}\mathbb{V}_0$ 的一个特征标。设 $\mathbb{V}_0(\chi)$ 是 $\mathbb{V}_0$ 对该特征标之核的商。于是有下述结果。

\begin{theorem}
    设 $\chi\in\operatorname{Op}_{{}^LG}(D)\subset\operatorname{Op}_{{}^LG}(D^\times)$。范畴 $\hat{\mathfrak{g}}_{\kappa_c}\text{-mod}^{G[[t]]}_\chi$ 等价于向量空间范畴：它在同构意义下唯一的不可约对象是 $\mathbb{V}_0(\chi)$，任何其他对象都同构于 $\mathbb{V}_0(\chi)$ 的若干拷贝的直和。
\end{theorem}

本定理可看作 10.3.1 节讨论的局部无分歧 Langlands 对应的范畴类比。它也为猜想 10.1.2 提供了第一个证据：我们看到，对一切 $\chi\in\operatorname{Op}_{{}^LG}(D)$，范畴 $\hat{\mathfrak{g}}_{\kappa_c}\text{-mod}^{G[[t]]}_\chi$ 彼此等价。

考虑整个正则 opers 族 $\operatorname{Op}^0_{{}^LG}=\operatorname{Op}_{{}^LG}(D)$ 而非单个正则 ${}^LG$-oper $\chi$，会更方便。设 $\hat{\mathfrak{g}}_{\kappa_c}\text{-mod}_{reg}$ 是范畴 $\hat{\mathfrak{g}}_{\kappa_c}\text{-mod}$ 的满子范畴，其对象是那些中心 $Z(\hat{\mathfrak{g}})$ 的作用通过同态
\[
    Z(\hat{\mathfrak{g}})\simeq\operatorname{Fun}\operatorname{Op}_{{}^LG}(D^\times)
    \to\operatorname{Fun}\operatorname{Op}_{{}^LG}(D)
\]
分解的 $\hat{\mathfrak{g}}_{\kappa_c}$-模。注意范畴 $\hat{\mathfrak{g}}_{\kappa_c}\text{-mod}_{reg}$ 是 10.1 节引入的范畴 $\hat{\mathfrak{g}}_{\kappa_c}\text{-mod}_Y$ 在 $Y=\operatorname{Op}_{{}^LG}(D)$ 时的例子。

设 $\hat{\mathfrak{g}}_{\kappa_c}\text{-mod}^{G[[t]]}_{reg}$ 是相应的 $G[[t]]$-等变范畴。把 $\hat{\mathfrak{g}}_{\kappa_c}\text{-mod}_{reg}$ 与 $\hat{\mathfrak{g}}_{\kappa_c}\text{-mod}^{G[[t]]}_{reg}$ 看作以 $\operatorname{Op}_{{}^LG}(D)$ 为底、纤维在 $\chi\in\operatorname{Op}_{{}^LG}(D)$ 上分别为 $\hat{\mathfrak{g}}_{\kappa_c}\text{-mod}_\chi$ 与 $\hat{\mathfrak{g}}_{\kappa_c}\text{-mod}^{G[[t]]}_\chi$ 的纤维范畴，是有启发性的。

现在我们来描述范畴 $\hat{\mathfrak{g}}_{\kappa_c}\text{-mod}^{G[[t]]}_{reg}$。这一描述特别蕴含定理 10.3.3。

为简化公式，下文对 $\operatorname{Fun}\operatorname{Op}_{{}^LG}(D)$ 采用如下记号：
\[
    \mathfrak{z}=\mathfrak{z}(\hat{\mathfrak{g}})=\operatorname{Fun}\operatorname{Op}_{{}^LG}(D).
\]
设 $\mathfrak{z}\text{-mod}$ 是交换代数 $\mathfrak{z}$ 上的模范畴。等价地，它是 $\operatorname{Op}_{{}^LG}(D)$ 上的拟凝聚层范畴。

按定义，$\hat{\mathfrak{g}}_{\kappa_c}\text{-mod}^{G[[t]]}_{reg}$ 的任何对象都是 $\mathfrak{z}$-模。引入函子
\[
    F:\hat{\mathfrak{g}}_{\kappa_c}\text{-mod}^{G[[t]]}_{reg}\to\mathfrak{z}\text{-mod},
    \qquad M\mapsto\operatorname{Hom}_{\hat{\mathfrak{g}}_{\kappa_c}}(\mathbb{V}_0,M);
\]
\[
    G:\mathfrak{z}\text{-mod}\to\hat{\mathfrak{g}}_{\kappa_c}\text{-mod}^{G[[t]]}_{reg},
    \qquad \mathcal{F}\mapsto\mathbb{V}_0\otimes_{\mathfrak{z}}\mathcal{F}.
\]
下述定理已在 [Frenkel and Gaitsgory 2004] 定理 6.3 中证明（此方向的重要结果此前已由 [Beilinson and Drinfeld 1997] 得到）。

\begin{theorem}
    函子 $F$ 与 $G$ 是互逆的范畴等价
    \begin{equation}
        \hat{\mathfrak{g}}_{\kappa_c}\text{-mod}^{G[[t]]}_{reg}\simeq\mathfrak{z}\text{-mod}.
        \tag{10.3-9}
    \end{equation}
\end{theorem}

我们将在下一小节给出证明。在此之前先指出，它立刻蕴含定理 10.3.3。事实上，对每个 $\chi\in\operatorname{Op}_{{}^LG}(D)$，范畴 $\hat{\mathfrak{g}}_{\kappa_c}\text{-mod}^{G[[t]]}_\chi$ 是 $\hat{\mathfrak{g}}_{\kappa_c}\text{-mod}^{G[[t]]}_{reg}$ 的满子范畴，其对象是那些作为 $\mathfrak{z}$-模被 $\chi$ 的极大理想 $I_\chi$ 零化的 $\hat{\mathfrak{g}}_{\kappa_c}$-模。由定理 10.3.4，此范畴等价于被 $I_\chi$ 零化的 $\mathfrak{z}$-模范畴。但后者是（概形论意义上）支撑在点 $\chi$ 处的 $\mathfrak{z}$-模范畴，它等价于向量空间范畴。

\subsection{定理 10.3.4 的证明}

首先注意，函子 $F$ 忠实，即若 $M\neq0$，则 $F(M)\neq0$。事实上，注意
\[
    F(M)=\operatorname{Hom}_{\hat{\mathfrak{g}}_{\kappa_c}}(\mathbb{V}_0,M)
    =M^{\mathfrak{g}[[t]]},
\]
因为任何 $\hat{\mathfrak{g}}_{\kappa_c}$-同态 $\mathbb{V}_0\to M$ 由 $\mathbb{V}_0$ 的生成向量在 $M$ 中的像唯一确定，而该生成向量是 $\mathfrak{g}[[t]]$-不变的。如引理 10.3.2 的证明所示，$\hat{\mathfrak{g}}_{\kappa_c}\text{-mod}^{G[[t]]}_{reg}$ 中任何非零 $\hat{\mathfrak{g}}_{\kappa_c}$-模 $M$ 必含有被 $\mathfrak{g}\otimes t\mathbb{C}[[t]]$ 零化的非零向量。由该向量生成的 $\mathfrak{g}$-模必是平凡的，否则中心在其上的作用不会通过 $\mathfrak{z}$ 分解。因此 $M$ 含有被 $\mathfrak{g}[[t]]$ 零化的非零向量，故 $M^{\mathfrak{g}[[t]]}\neq0$。

其次，我们注意到函子 $G$ 是函子 $F$ 的左伴随，即有一族相容的同构
\begin{equation}
    \operatorname{Hom}(\mathcal{F},F(M))\simeq\operatorname{Hom}(G(\mathcal{F}),M).
    \tag{10.3-10}
\end{equation}
特别地，在此式中取 $M=G(\mathcal{F})$，得到一族相容的映射 $\mathcal{F}\to F\circ G(\mathcal{F})$，即
\begin{equation}
    \mathcal{F}\to\operatorname{Hom}_{\hat{\mathfrak{g}}_{\kappa_c}}
    \Bigl(\mathbb{V}_0,\mathbb{V}_0\otimes_{\mathfrak{z}}\mathcal{F}\Bigr).
    \tag{10.3-11}
\end{equation}
我们断言这实际上是同构。为此用到下面两个结果。

\begin{theorem}
    $\mathfrak{z}$-模 $\mathbb{V}_0$ 是自由的。
\end{theorem}

\begin{proof}
    回忆 $\mathbb{V}_0$ 上有 PBW 滤过，其分次伴随同构于 $\operatorname{Fun}\mathfrak{g}^*[[t]]$（见 3.3.3 节）。PBW 滤过在 $\mathfrak{z}(\hat{\mathfrak{g}})=(\mathbb{V}_0)^{\mathfrak{g}[[t]]}$ 上诱导滤过，其分次伴随同构于 $\operatorname{Inv}\mathfrak{g}^*[[t]]=(\operatorname{Fun}\mathfrak{g}^*[[t]])^{\mathfrak{g}[[t]]}$（见命题 4.3.3）。据 [Mustată 2001]，$\operatorname{Fun}\mathfrak{g}^*[[t]]$ 是 $\operatorname{Inv}\mathfrak{g}^*[[t]]$ 上的自由模。由此直接推出原滤过对象也有同样性质，故 $\mathbb{V}_0$ 是 $\mathfrak{z}=\operatorname{Fun}\operatorname{Op}_{{}^LG}(D)$ 上的自由模。
\end{proof}

% [Tutor Note: “由此直接推出原滤过对象也有同样性质”省略的是 Rees 模论证：若分次伴随 $\operatorname{gr}\mathbb{V}_0$ 在 $\operatorname{gr}\mathfrak{z}$ 上自由，则对滤过有界的模可取生成元在分次中的提升，用归纳证明原模自由。此处需要 $\mathfrak{z}(\hat{\mathfrak{g}})=(\mathbb{V}_0)^{\mathfrak{g}[[t]]}$ 与命题 4.3.3 的分次等同。]

这立即给出

\begin{corollary}
    函子 $G$ 正合。
\end{corollary}

其次，设 $\hat{\mathfrak{g}}_{\kappa_c}\text{-mod}^{G[[t]]}$ 是 $(\hat{\mathfrak{g}}_{\kappa_c},G[[t]])$ Harish--Chandra 模范畴（对中心的作用不加任何限制）。设
\[
    \operatorname{Ext}^i_{\hat{\mathfrak{g}},G[[t]]}(\mathbb{V}_0,M)
\]
是函子 $M\mapsto\operatorname{Hom}_{\hat{\mathfrak{g}}_{\kappa_c}}(\mathbb{V}_0,M)$ 从 $\hat{\mathfrak{g}}_{\kappa_c}\text{-mod}^{G[[t]]}$ 到向量空间范畴的高阶导出函子。

\begin{theorem}[Frenkel and Teleman 2006]
    有
    \begin{equation}
        \operatorname{Ext}^i_{\hat{\mathfrak{g}},G[[t]]}(\mathbb{V}_0,\mathbb{V}_0)
        \simeq\Omega^i_{\mathfrak{z}},
        \tag{10.3-12}
    \end{equation}
    即 $\operatorname{Op}_{{}^LG}(D)$ 上 $i$ 次微分形式构成的空间。
\end{theorem}

注意这是同构 $\operatorname{Hom}_{\hat{\mathfrak{g}}_{\kappa_c}}(\mathbb{V}_0,\mathbb{V}_0)\simeq\mathfrak{z}$ 的推广，后者对应公式 (10.3-12) 中 $i=0$ 的情形。\transnote{这个同构的几何来源是：$\operatorname{Ext}^i_{\hat{\mathfrak{g}},G[[t]]}(\mathbb{V}_0,\mathbb{V}_0)$ 可实现为 Chevalley 复形 $C^\bullet(\mathfrak{g}[[t]],\mathfrak{g};\mathbb{V}_0)$ 的上同调，而 Drinfeld--Sokolov 约化把该上同调认同为 $\operatorname{Op}_{{}^LG}(D)$ 上的 $i$-形式。等价地，$\Omega^i_{\mathfrak{z}}$ 是 $\mathfrak{z}=\operatorname{Fun}\operatorname{Op}_{{}^LG}(D)$ 的 $i$-次 Kähler 微分形式；右端因此自动带有 $\mathfrak{z}$-模结构，这正是下面能对 $\mathcal{F}$ 用谱序列的依据。$i=0$ 时它退化为 $\mathfrak{z}$ 本身，即 (9.5-1) 的中心同构。[来源：原书 p.323；E. Frenkel and C. Teleman, Self-extensions of Verma modules and differential forms on opers, Compos. Math. 142 (2006), 477--500.]}

% [Tutor Note: 小心上同调的下标：定理 10.3.7 中记作 $\operatorname{Ext}^i_{\hat{\mathfrak{g}},G[[t]]}$（水平未固定），而 (10.3-13) 中为 $\operatorname{Ext}^i_{\hat{\mathfrak{g}}_{\kappa_c},G[[t]]}$。两者通过中心 $Z(\hat{\mathfrak{g}})$ 的特征标相联系；把 $\hat{\mathfrak{g}}$ 与 $\hat{\mathfrak{g}}_{\kappa_c}$ 混淆会直接导致 (10.3-13) 的证明失真。]

现在我们来证明如下同构：
\begin{equation}
    \operatorname{Ext}^i_{\hat{\mathfrak{g}}_{\kappa_c},G[[t]]}(\mathbb{V}_0,\mathbb{V}_0)
    \otimes_{\mathfrak{z}}\mathcal{F}
    \simeq
    \operatorname{Ext}^i_{\hat{\mathfrak{g}}_{\kappa_c},G[[t]]}
    \Bigl(\mathbb{V}_0,\mathbb{V}_0\otimes_{\mathfrak{z}}\mathcal{F}\Bigr),
    \qquad i\ge0.
    \tag{10.3-13}
\end{equation}
当 $i=0$ 时即得到映射 (10.3-11) 是同构，这正是所需。

为证同构 (10.3-13)，注意函子
\[
    \mathcal{F}\mapsto\operatorname{Ext}^i_{\hat{\mathfrak{g}}_{\kappa_c},G[[t]]}
    \Bigl(\mathbb{V}_0,\mathbb{V}_0\otimes_{\mathfrak{z}}\mathcal{F}\Bigr)
\]
与取正向极限可交换。这可由它们实现为计算相对李代数上同调 $H^i\bigl(\mathfrak{g}[[t]],\mathfrak{g},\mathbb{V}_0\otimes_{\mathfrak{z}}\mathcal{F}\bigr)$ 的 Chevalley 复形的上同调看出，见 [Frenkel and Teleman 2006] 命题 2.1 与引理 3.1。因此不失一般性，可以也的确假设 $\mathcal{F}$ 是有限表现的 $\mathfrak{z}$-模。

由于 $\mathfrak{z}$ 同构于自由多项式代数（见公式 (4.3-1)），任何这样的模都有有限投射消解
\[
    0\to\mathcal{P}_n\to\cdots\to\mathcal{P}_1\to\mathcal{P}_0\to\mathcal{F}\to0.
\]
按同调代数的标准结果，(10.3-13) 的右端可由谱序列计算，其 $E^{p,q}_1$ 项同构于
\[
    \operatorname{Ext}^p_{\hat{\mathfrak{g}}_{\kappa_c},G[[t]]}
    \Bigl(\mathbb{V}_0,\mathbb{V}_0\otimes_{\mathfrak{z}}\mathcal{P}_q\Bigr).
\]
由于每个 $\mathcal{P}_j$ 投射，有
\[
    \operatorname{Ext}^p_{\hat{\mathfrak{g}}_{\kappa_c},G[[t]]}
    \Bigl(\mathbb{V}_0,\mathbb{V}_0\otimes_{\mathfrak{z}}\mathcal{P}_q\Bigr)
    =\operatorname{Ext}^p_{\hat{\mathfrak{g}}_{\kappa_c},G[[t]]}(\mathbb{V}_0,\mathbb{V}_0)
    \otimes_{\mathfrak{z}}\mathcal{P}_q.
\]
其次，由定理 10.3.7，每个 $\operatorname{Ext}^p_{\hat{\mathfrak{g}}_{\kappa_c},G[[t]]}(\mathbb{V}_0,\mathbb{V}_0)$ 是自由 $\mathfrak{z}$-模。因此谱序列的第二项等于
\[
    E^{p,0}_2=\operatorname{Ext}^p_{\hat{\mathfrak{g}}_{\kappa_c},G[[t]]}(\mathbb{V}_0,\mathbb{V}_0)
    \otimes_{\mathfrak{z}}\mathcal{F},
\]
且 $q\neq0$ 时 $E^{p,q}=0$。故谱序列在第二项退化，结果就是同构 (10.3-13)。

于是我们证明了 $F\circ G\simeq\operatorname{Id}$。下面证明 $G\circ F\simeq\operatorname{Id}$。

注意由伴随 (10.3-10) 有映射 $G\circ F\to\operatorname{Id}$，即对 $M\in\hat{\mathfrak{g}}_{\kappa_c}\text{-mod}^{G[[t]]}_{reg}$ 有相容的一族映射
\begin{equation}
    G\circ F(M)\to M.
    \tag{10.3-14}
\end{equation}
我们需要证明此映射是同构。

先证 (10.3-14) 是单射。若 $M'$ 是 $G\circ F(M)\to M$ 的核，则由 $F$ 左正合可得
\[
    F(M')=\operatorname{Ker}\bigl(F\circ G\circ F(M)\to F(M)\bigr)
    \simeq\operatorname{Ker}\bigl(F(M)\to F(M)\bigr)=0.
\]
但已知 $F$ 忠实，故 $M'=0$。

其次证 (10.3-14) 是满射。设 $M''$ 是 $G\circ F(M)\to M$ 的余核。有长正合列
\begin{equation}
    0\to F\circ G\circ F(M)\to F(M)\to F(M'')\to R^1F(G\circ F(M))\to\cdots,
    \tag{10.3-15}
\end{equation}
其中 $R^1F$ 是函子 $F$ 的第一右导出函子。因此
\[
    R^1F(L)=\operatorname{Ext}^1_{\hat{\mathfrak{g}}_{\kappa_c}\text{-mod}^{G[[t]]}_{reg}}
    (\mathbb{V}_0,L).
\]
为完成证明，需要下述结果：

\begin{lemma}
    对任何 $\mathfrak{z}$-模 $\mathcal{F}$，有
    \[
        \operatorname{Ext}^1_{\hat{\mathfrak{g}}_{\kappa_c}\text{-mod}^{G[[t]]}_{reg}}
        \Bigl(\mathbb{V}_0,\mathbb{V}_0\otimes_{\mathfrak{z}}\mathcal{F}\Bigr)=0.
    \]
\end{lemma}

\begin{proof}
    当 $\mathcal{F}=\mathfrak{z}$ 时，这已在 [Beilinson and Drinfeld 1997] 中证明。这里我们按 [Frenkel and Gaitsgory 2004] 对一般情形给出的证明。

    $\operatorname{Ext}^1_{\hat{\mathfrak{g}}_{\kappa_c}\text{-mod}^{G[[t]]}_{reg}}\bigl(\mathbb{V}_0,\mathbb{V}_0\otimes_{\mathfrak{z}}\mathcal{F}\bigr)=\operatorname{Ext}^1_{\hat{\mathfrak{g}}_{\kappa_c},G[[t]]}\bigl(\mathbb{V}_0,\mathbb{V}_0\otimes_{\mathfrak{z}}\mathcal{F}\bigr)$ 的任何元素都定义了 $\hat{\mathfrak{g}}_{\kappa_c}$-模的扩张
    \begin{equation}
        0\to\mathbb{V}_0\otimes_{\mathfrak{z}}\mathcal{F}\to\mathcal{M}\to\mathbb{V}_0\to0.
        \tag{10.3-16}
    \end{equation}
    于是 $\mathcal{M}$ 是范畴 $\hat{\mathfrak{g}}_{\kappa_c}\text{-mod}^{G[[t]]}_{reg}$ 的两个对象的扩张，中心 $Z(\hat{\mathfrak{g}})$ 通过商 $Z(\hat{\mathfrak{g}})\to\mathfrak{z}$ 作用。但模 $\mathcal{M}$ 未必是 $\hat{\mathfrak{g}}_{\kappa_c}\text{-mod}^{G[[t]]}_{reg}$ 的对象。命题的陈述意味着：若 $\mathcal{M}$ 是 $\hat{\mathfrak{g}}_{\kappa_c}\text{-mod}^{G[[t]]}_{reg}$ 的对象，则它作为 $\hat{\mathfrak{g}}_{\kappa_c}$-模必分裂。

    设 $I$ 是 $Z(\hat{\mathfrak{g}})$ 中由 $\mathfrak{z}(\hat{\mathfrak{g}})$ 生成的理想。\transnote{准确地说，$I$ 是由子代数 $\mathfrak{z}(\hat{\mathfrak{g}})=\operatorname{Fun}\operatorname{Op}_{{}^LG}(D)$ 在 $Z(\hat{\mathfrak{g}})=\operatorname{Fun}\operatorname{Op}_{{}^LG}(D^\times)$ 中生成的理想，即限制映射 $Z(\hat{\mathfrak{g}})\to\mathfrak{z}(\hat{\mathfrak{g}})$ 的核；它并不是 $\mathfrak{z}(\hat{\mathfrak{g}})$ 内部的理想。$\mathcal{M}$ 属于 $\hat{\mathfrak{g}}_{\kappa_c}\text{-mod}^{G[[t]]}_{reg}$ 当且仅当中心作用通过 $\mathfrak{z}(\hat{\mathfrak{g}})$ 分解，也就是 $I$ 在 $\mathcal{M}$ 上作用为零。于是 $I/I^2$ 可看作嵌入 $\operatorname{Op}(D)\subset\operatorname{Op}(D^\times)$ 的余法模，后面的锚映射 (10.3-19) 正来自此。[来源：原书 p.325；A. Beilinson and V. Drinfeld, Quantization of Hitchin's integrable system and Hecke eigensheaves, 1997, §3.6.]}则 $\mathcal{M}$ 是 $\hat{\mathfrak{g}}_{\kappa_c}\text{-mod}^{G[[t]]}_{reg}$ 的对象，当且仅当 $I$ 在它上面作用为 $0$。注意 $I$ 必在子模 $\mathbb{V}_0\otimes_{\mathfrak{z}}\mathcal{F}$ 与商 $\mathbb{V}_0$ 上为零。因此，取 (10.3-16) 的线性分裂并对 $\mathbb{V}_0$ 在 $\mathcal{M}$ 中的像作用 $I$，得到映射
    \[
        I\to\operatorname{Hom}_{\hat{\mathfrak{g}}_{\kappa_c}}
        \Bigl(\mathbb{V}_0,\mathbb{V}_0\otimes_{\mathfrak{z}}\mathcal{F}\Bigr).
    \]
    而且 $I^2\subset I$ 的任何元素都映到 $0$。于是得到映射
    \[
        (I/I^2)\otimes_{\mathfrak{z}}\operatorname{Ext}^1_{\hat{\mathfrak{g}}_{\kappa_c},G[[t]]}
        \Bigl(\mathbb{V}_0,\mathbb{V}_0\otimes_{\mathfrak{z}}\mathcal{F}\Bigr)
        \to\operatorname{Hom}_{\hat{\mathfrak{g}}_{\kappa_c}}
        \Bigl(\mathbb{V}_0,\mathbb{V}_0\otimes_{\mathfrak{z}}\mathcal{F}\Bigr).
    \]
    我们需要证明：对 $\operatorname{Ext}^1_{\hat{\mathfrak{g}}_{\kappa_c},G[[t]]}\bigl(\mathbb{V}_0,\mathbb{V}_0\otimes_{\mathfrak{z}}\mathcal{F}\bigr)$ 的任何元素，相应的映射 $I/I^2\to\operatorname{Hom}_{\hat{\mathfrak{g}}_{\kappa_c}}\bigl(\mathbb{V}_0,\mathbb{V}_0\otimes_{\mathfrak{z}}\mathcal{F}\bigr)$ 非零。等价地，需证映射
    \[
        \operatorname{Ext}^1_{\hat{\mathfrak{g}}_{\kappa_c},G[[t]]}
        \Bigl(\mathbb{V}_0,\mathbb{V}_0\otimes_{\mathfrak{z}}\mathcal{F}\Bigr)
        \to\operatorname{Hom}_{\mathfrak{z}}\Bigl(I/I^2,
        \operatorname{Hom}_{\hat{\mathfrak{g}}_{\kappa_c}}
        \bigl(\mathbb{V}_0,\mathbb{V}_0\otimes_{\mathfrak{z}}\mathcal{F}\bigr)\Bigr)
    \]
    的核为 $0$。事实上，这个核等于 $\operatorname{Ext}^1_{\hat{\mathfrak{g}}_{\kappa_c}\text{-mod}^{G[[t]]}_{reg}}\bigl(\mathbb{V}_0,\mathbb{V}_0\otimes_{\mathfrak{z}}\mathcal{F}\bigr)$，其消失正是我们要证的。

    现在由定理 10.3.7 与公式 (10.3-13)，最后一个映射可改写为如下形式：
    \begin{equation}
        \Omega^1_{\mathfrak{z}}\otimes_{\mathfrak{z}}\mathcal{F}
        \to\operatorname{Hom}_{\mathfrak{z}}\Bigl(I/I^2,\mathcal{F}\Bigr).
        \tag{10.3-17}
    \end{equation}
    为写出此映射的显式公式，我们按 [Frenkel and Teleman 2006] 说明如何构造同构
    \begin{equation}
        \Omega^1_{\mathfrak{z}}\otimes_{\mathfrak{z}}\mathcal{F}
        \simeq\operatorname{Ext}^1_{\hat{\mathfrak{g}}_{\kappa_c},G[[t]]}
        \Bigl(\mathbb{V}_0,\mathbb{V}_0\otimes_{\mathfrak{z}}\mathcal{F}\Bigr),
        \tag{10.3-18}
    \end{equation}
    换言之，如何从 $\Omega^1_{\mathfrak{z}}\otimes_{\mathfrak{z}}\mathcal{F}$ 的元素出发构造形如 (10.3-16) 的扩张 $\mathcal{M}$。为简单起见设 $\mathcal{F}=\mathfrak{z}$（一般情形类似）。则作为向量空间，扩张 $\mathcal{M}$ 是 $\mathbb{V}_0$ 的两个拷贝的直和。由线性性，可假设 $\Omega^1_{\mathfrak{z}}$ 中我们的元素形如 $B\,dA$，其中 $A,B\in\mathfrak{z}$。由于 $\mathbb{V}_0$ 由真空向量 $|0\rangle$ 生成，扩张 $\mathcal{M}$ 由 $\mathfrak{g}[[t]]$ 在对应于商 $\mathbb{V}_0$ 的直和项中向量 $|0\rangle$ 上的作用唯一确定。于是令
    \[
        X\cdot|0\rangle=B\cdot\lim_{\epsilon\to0}\frac{1}{\epsilon}X\cdot A_\epsilon,
        \qquad X\in\mathfrak{g}[[t]].
    \]
    这里 $A_\epsilon$ 是 $A\in\mathfrak{z}(\hat{\mathfrak{g}})\subset\mathbb{V}_0=V_{\kappa_c}(\mathfrak{g})$ 的任意形变，看作水平为 $\kappa_c+\epsilon\kappa_0$ 的真空模 $V_{\kappa_c+\epsilon\kappa_0}(\mathfrak{g})$ 的元素（其中 $\kappa_0$ 是 $\mathfrak{g}$ 上非零的不变内积）。由于在 $\mathbb{V}_0=V_{\kappa_c}(\mathfrak{g})$ 中对一切 $X\in\mathfrak{g}[[t]]$ 有 $X\cdot A=0$，此公式的右端良定义。

    现在回忆中心 $Z(\hat{\mathfrak{g}})$ 是 Poisson 代数，其 Poisson 括号在 8.3.1 节定义。此外，由定理 8.3.1 与引理 8.3.2，此 Poisson 代数同构于 Poisson 代数
    \[
        \operatorname{Fun}\operatorname{Op}_{{}^LG}(D^\times)
    \]
    配上经 Drinfeld--Sokolov 约化得到的 Poisson 代数结构。容易看出，关于此 Poisson 结构，理想 $I$ 是 Poisson 的，即 $\{I,I\}\subset I$。在此情形，Poisson 括号映射
    \[
        \operatorname{Fun}\operatorname{Op}_{{}^LG}(D^\times)
        \to\Theta\bigl(\operatorname{Op}_{{}^LG}(D^\times)\bigr),
    \]
    其中 $\Theta(\operatorname{Op}_{{}^LG}(D^\times))$ 是 $\operatorname{Op}_{{}^LG}(D^\times)$ 上向量场的空间，诱导出映射
    \begin{equation}
        I/I^2\to\Theta_{\mathfrak{z}},
        \tag{10.3-19}
    \end{equation}
    其中 $\Theta_{\mathfrak{z}}$ 是 $\operatorname{Op}_{{}^LG}(D)$ 上向量场的空间。

    由上述同构 (10.3-18) 的描述与 8.3.1 节中 $Z(\hat{\mathfrak{g}})$ 的 Poisson 结构的定义可知，映射 (10.3-17) 是重言同构
    \[
        \Omega^1_{\mathfrak{z}}\otimes_{\mathfrak{z}}\mathcal{F}
        \xrightarrow{\ \sim\ }\operatorname{Hom}_{\mathfrak{z}}(\Theta_{\mathfrak{z}},\mathcal{F})
    \]
    （注意 $\Omega^1_{\mathfrak{z}}$ 是自由 $\mathfrak{z}$-模）与由 (10.3-19) 诱导的映射
    \begin{equation}
        \operatorname{Hom}_{\mathfrak{z}}(\Theta_{\mathfrak{z}},\mathcal{F})
        \to\operatorname{Hom}_{\mathfrak{z}}\Bigl(I/I^2,\mathcal{F}\Bigr)
        \tag{10.3-20}
    \end{equation}
    的复合。但据 [Beilinson and Drinfeld 1997] 定理 3.6.7，映射 (10.3-19) 是满的（事实上 $I/I^2$ 同构于 $\operatorname{Op}_{{}^LG}(D)$ 上对应于万有 ${}^LG$-丛的李代数胚，而 (10.3-19) 是对应的锚映射，故其核作为 $\mathfrak{z}$-模自由）。因此对任何 $\mathfrak{z}$-模 $\mathcal{F}$，映射 (10.3-20) 是单射。从而映射 (10.3-17) 也是单射。证毕。
\end{proof}

由引理 10.3.8，序列 (10.3-15) 给出短正合列
\[
    0\to F\circ G\circ F(M)\to F(M)\to F(M'')\to0.
\]
但第一个箭头是同构，因为 $F\circ G\simeq\operatorname{Id}$，故 $F(M'')=0$，从而 $M''=0$。

定理 10.3.4、从而定理 10.3.3，至此证毕。

\subsection{球面 Hecke 代数的作用}

在 10.3.1 节我们讨论了群 $G(F)$ 的不可约无分歧表示，其中 $F$ 是局部非阿基米德域。我们看到，这样的表示由 Langlands 对偶群 ${}^LG$ 的共轭类参数化。给定这样的共轭类 $\gamma$，有不可约无分歧表示 $(R_\gamma,\pi_\gamma)$，它含有一维子空间 $(R_\gamma)^{\pi_\gamma(K_0)}$，由 $K_0$-不变向量组成。球面 Hecke 代数 $H(G(F),K_0)$ 经 Satake 同构同构于 $\operatorname{Rep}{}^LG$，它在此空间上按特征标 $\phi_\gamma$ 作用，见公式 (10.3-5)。

在几何情形，我们已经论证：对任何 $\chi\in\operatorname{Op}_{{}^LG}(D)$，配备圈群 $G((t))$ 作用的范畴 $\hat{\mathfrak{g}}_{\kappa_c}\text{-mod}_\chi$ 应被看作 $(R_\gamma,\pi_\gamma)$ 的范畴化。此外，它的 $(\hat{\mathfrak{g}}_{\kappa_c},G[[t]])$ Harish--Chandra 模子范畴 $\hat{\mathfrak{g}}_{\kappa_c}\text{-mod}^{G[[t]]}_\chi$ 应被看作一维空间 $(R_\gamma)^{\pi_\gamma(K_0)}$ 的“范畴化”。由定理 10.3.3，后一范畴等价于向量空间范畴，而向量空间范畴确实是一维向量空间的范畴化。因此此结果与经典图像一致。

现在我们讨论球面 Hecke 代数在这个一维空间上的作用的范畴类比。

如 10.2.3 节所述，球面 Hecke 代数的范畴类比是仿射 Grassmann 流形 $\operatorname{Gr}=G((t))/G[[t]]$ 上 $G[[t]]$-等变 $D$-模构成的范畴。$\operatorname{Gr}$ 与此范畴的精确定义见 [Beilinson and Drinfeld 1997; Frenkel and Gaitsgory 2006c]。无分歧情形有一个重要性质：与这些 $D$-模的卷积函子正合，这意味着我们无需考虑导出范畴；这样的 $D$-模的阿贝尔范畴就够用了。记此阿贝尔范畴为 $\mathcal{H}(G((t)),G[[t]])$。这是球面 Hecke 代数 $H(G((t)),G[[t]])$ 的范畴版本。\transnote{球面 Hecke 范畴的对象是 $\operatorname{Gr}$ 上右 $G[[t]]$-等变的 $D$-模；其张量积是 Hecke 卷积（沿群作用图定义），单位对象是原点处的 $\delta$-模。这里“球面”与经典情形同义：只保留 $G[[t]]$-等变的数据；若改用 Iwahori 子群 $I$，得到的是定义在仿射旗簇 $G((t))/I$ 上的 Iwahori Hecke 范畴，其对应的 Iwahori Hecke 代数不再交换，卷积也不逐项正合，通常要放在导出范畴或混合范畴中处理。无分歧情形卷积函子正合，故本书只需阿贝尔范畴——这是球面情形特有的简化。[来源：原书 p.327--328；A. Beilinson and V. Drinfeld, Quantization of Hitchin's integrable system and Hecke eigensheaves, 1997, §4; E. Frenkel and D. Gaitsgory, Local geometric Langlands correspondence and affine Kac-Moody algebras, 2006, §2.]}

由 [Mirković and Vilonen 2007] 的结果，范畴 $\mathcal{H}(G((t)),G[[t]])$ 带有自然的张量范畴结构，它等价于 ${}^LG$ 的表示张量范畴 $\operatorname{Rep}{}^LG$。这应被看作 Satake 同构的范畴类比。\transnote{“范畴化”在这里有确切含义：Satake 同构 (10.3-3) 只是交换代数 $H(G(F),K_0)$ 与 $\operatorname{Rep}{}^LG$ 的同构，而 Mirković--Vilonen 定理把 $\mathcal{H}(G((t)),G[[t]])$ 与 $\operatorname{Rep}{}^LG$ 作成张量范畴的等价——两个范畴的对象、态射、张量积与单位对象逐一对应，不只是 Grothendieck 群相同。于是“$H$-模的本征值”被提升为“$\mathcal{H}$-模的本征对象”，(10.3-21) 中的 $\alpha_V$ 正是这一提升的显式形式。也正因如此，Hecke 本征模条件 (10.3-26) 才能被读作 (10.3-5) 的范畴类比。[来源：原书 p.328；I. Mirković and K. Vilonen, Geometric Langlands duality and representations of algebraic groups over commutative rings, Ann. of Math. 166 (2007), 95--143.]}于是对 $\operatorname{Rep}{}^LG$ 的每个对象 $V$，有 $\mathcal{H}(G((t)),G[[t]])$ 的一个对象，记作 $\mathcal{H}_V$。Hecke 本征向量性质 (10.3-5) 的类比应当是什么？

如 10.2.3 节所述，范畴 $\mathcal{H}(G((t)),G[[t]])$ 自然地作用在范畴 $\hat{\mathfrak{g}}_{\kappa_c}\text{-mod}^{G[[t]]}_\chi$ 上，此作用应被看作 $H(G(F),K_0)$ 在 $(R_\gamma)^{\pi_\gamma(K_0)}$ 上作用的范畴类比。

现在由定理 10.3.3，$\hat{\mathfrak{g}}_{\kappa_c}\text{-mod}^{G[[t]]}_\chi$ 的任何对象都是 $\mathbb{V}_0(\chi)$ 的若干拷贝的直和。因此只需描述 $\mathcal{H}(G((t)),G[[t]])$ 在 $\mathbb{V}_0(\chi)$ 上的作用。此作用由下述陈述描述（它来自 [Beilinson and Drinfeld 1997]）：存在一族同构
\begin{equation}
    \alpha_V:\mathcal{H}_V\star\mathbb{V}_0(\chi)\xrightarrow{\ \sim\ }
    \underline{V}\otimes\mathbb{V}_0(\chi),
    \qquad V\in\operatorname{Rep}{}^LG,
    \tag{10.3-21}
\end{equation}
其中 $\underline{V}$ 是表示 $V$ 底层的向量空间（细节见 [Frenkel and Gaitsgory 2006c; Frenkel and Gaitsgory 2005]）。此外，这些同构与 $\mathcal{H}_V$ 上的张量积结构（由卷积给出）及 $\underline{V}$ 上的张量积结构（由向量空间张量积给出）相容。

% [Tutor Note: 式 (10.3-21) 中 $\underline{V}$ 的下划线表示“$V$ 的底向量空间”，即忘掉 ${}^LG$-作用后的向量空间；这一区分在本小节很关键，因为 Hecke 本征模条件把标量乘法换成了“乘以向量空间”。]

鉴于定理 10.3.3，这并不奇怪。事实上，由定义，$\mathcal{H}_V\star\mathbb{V}_0(\chi)$ 仍是范畴 $\hat{\mathfrak{g}}_{\kappa_c}\text{-mod}^{G[[t]]}_\chi$ 的对象。因此它必同构于 $U_V\otimes_{\mathbb{C}}\mathbb{V}_0(\chi)$，其中 $U_V$ 是向量空间。但这样我们得到函子
\[
    \mathcal{H}(G((t)),G[[t]])\to\operatorname{Vect},
    \qquad \mathcal{H}_V\mapsto U_V.
\]
由构造可知这是张量函子。因此标准的 Tannaka 形式体系蕴含 $U_V$ 同构于 $\underline{V}$。

同构 (10.3-21) 应被看作 Hecke 本征向量条件 (10.3-5) 的范畴类比。区别在于：(10.3-5) 中 Hecke 代数元素在 $R_\gamma$ 的 $K_0$-不变向量上的作用归结为乘以一个标量；而 Hecke 范畴 $\mathcal{H}(G((t)),G[[t]])$ 的对象在 $\hat{\mathfrak{g}}_{\kappa_c}\text{-mod}_\chi$ 的 $G[[t]]$-等变对象 $\mathbb{V}_0(\chi)$ 上的作用，归结为乘以一个向量空间，即 ${}^LG$ 的相应表示底层的向量空间。满足此性质的模自然地称为\textbf{Hecke 本征模}。于是我们得到 $\mathbb{V}_0(\chi)$ 是 Hecke 本征模。这与我们的预期一致：范畴 $\hat{\mathfrak{g}}_{\kappa_c}\text{-mod}^{G[[t]]}_\chi$ 是 $R_\gamma$ 中 $K_0$-不变向量空间的范畴版本。

几何情形中缺失的一个要素是 ${}^LG$ 的共轭类 $\gamma$。我们回忆，在经典 Langlands 对应中它曾是伽罗瓦群 $\operatorname{Gal}(\overline{\mathbb{F}}_q/\mathbb{F}_q)$ 的 Frobenius 元素的像；但在几何情形（基域为代数闭域 $\mathbb{C}$）中没有类似物。因此，虽然经典情形中无分歧局部系统由共轭类 $\gamma$ 参数化，几何情形中在同构意义下只有一个无分歧局部系统。然而这个局部系统有很大的自同构群，即 ${}^LG$ 本身。我们将论证，几何情形中替代 $\gamma$ 的，是群 ${}^LG$ 通过范畴 $\hat{\mathfrak{g}}_{\kappa_c}\text{-mod}_\chi$ 的自动同构所起的作用。我们将在下两小节讨论这一点。\transnote{经典情形中 $\gamma$ 是 Frobenius 的像，给出 ${}^LG$ 中一个具体的半单共轭类；几何情形没有 Frobenius，能保留的只是 Hecke 本征模数据——对每个 $V\in\operatorname{Rep}{}^LG$，一族与张量积相容的同构 $\alpha_V$。把 oper 底层的 ${}^LG$-torsor $\mathcal{F}_0$ 的所有平凡化放在一起，便得到下一小节的群胚 $\operatorname{Isom}$；$\gamma$ 的共轭类被 $\operatorname{Isom}$ 在纤维范畴上的作用取代。这不是把 $\gamma$ 丢掉，而是把它范畴化为自同构的相干数据。[来源：原书 p.329、p.334；E. Frenkel and D. Gaitsgory, Local geometric Langlands correspondence and affine Kac-Moody algebras, 2006, §2.]}

% [Tutor Note: 这一段是理解 10.3.7 节群胚作用的钥匙：几何情形没有 Frobenius，${}^LG$ 的“共轭类”信息只能由范畴的自同构（或等价的 Hecke 本征模数据）承载，因此 $\gamma$ 被 ${}^LG$ 在范畴上的作用取代。注意这里的 ${}^LG$ 是紧李群的复化，其作用不保持范畴中的具体模，而只保持范畴的等价类结构。]

\subsection{表示范畴与 $D$-模范畴}

当我们在 1.3.4 节讨论表示的范畴化程序时，看到构造配备圈群 $G((t))$ 作用的范畴有两种可能的方案。第一种考虑 ind-概形 $G((t))/K$ 上 $D$-模构成的范畴，其中 $K$ 是 $G((t))$ 的“紧”子群，例如 $G[[t]]$ 或 Iwahori 子群。第二种考虑表示范畴 $\hat{\mathfrak{g}}_{\kappa_c}\text{-mod}_\chi$。到目前为止我们只专注于第二种方案，但讨论第一类范畴也是有启发性的。

在 1.3.3 节考虑的玩具模型中，我们讨论了固定中心特征标的 $\mathfrak{g}$-模范畴以及旗簇 $G/B$ 上的 $D$-模范畴。我们论证过，两者都可看作群 $G(\mathbb{F}_q)$ 在 $(G/B)(\mathbb{F}_q)$ 上的函数空间中的表示的范畴化。按 Beilinson--Bernstein 理论，这两个范畴等价，全局截面函子连接二者。对仿射 Kac--Moody 代数，类似的事情是否也成立？

仿射 Grassmann 流形 $\operatorname{Gr}=G((t))/G[[t]]$ 可看作圈群 $G((t))$ 的旗簇 $G/B$ 的最简单类比。考虑 $G((t))/G[[t]]$ 上的 $D$-模范畴（精确定义见 [Beilinson and Drinfeld 1997; Frenkel and Gaitsgory 2006c]）。有从此范畴到 $\mathfrak{g}((t))$-模范畴的全局截面函子。为得到 $\hat{\mathfrak{g}}_{\kappa_c}$-模，我们需要改用由线丛 $\mathcal{L}_{\kappa_c}$ 扭变的 $D$-模范畴 $\mathcal{D}_{\kappa_c}\text{-mod}$。这是 $\operatorname{Gr}$ 上唯一一个带有 $\hat{\mathfrak{g}}_{\kappa_c}$ 的作用（中心元 $\mathbf{1}$ 映到恒等）并提升 $\mathfrak{g}((t))$ 在 $\operatorname{Gr}$ 上自然作用的线丛 $\mathcal{L}_{\kappa_c}$。于是对 $\mathcal{D}_{\kappa_c}\text{-mod}$ 的任何对象 $\mathcal{M}$，全局截面空间 $\Gamma(\operatorname{Gr},\mathcal{M})$ 是 $\hat{\mathfrak{g}}_{\kappa_c}$-模。此外已知（见 [Beilinson and Drinfeld 1997; Frenkel and Gaitsgory 2004]），$\Gamma(\operatorname{Gr},\mathcal{M})$ 实际上是 $\hat{\mathfrak{g}}_{\kappa_c}\text{-mod}_{reg}$ 的对象。因此有全局截面函子
\[
    \Gamma:\mathcal{D}_{\kappa_c}\text{-mod}\to\hat{\mathfrak{g}}_{\kappa_c}\text{-mod}_{reg}.
\]
我们注意，范畴 $\mathcal{D}\text{-mod}$ 与 $\mathcal{D}_{\kappa_c}\text{-mod}$ 通过函子 $\mathcal{M}\mapsto\mathcal{M}\otimes\mathcal{L}_{\kappa_c}$ 等价。但相应的全局截面函子差别很大。

然而，与 Beilinson--Bernstein 情形不同，函子 $\Gamma$ 绝不可能是范畴等价。原因有二。首先，范畴 $\hat{\mathfrak{g}}_{\kappa_c}\text{-mod}_{reg}$ 有很大的中心，即代数 $\mathfrak{z}=\operatorname{Fun}\operatorname{Op}_{{}^LG}(D)$，而范畴 $\mathcal{D}_{\kappa_c}\text{-mod}$ 的中心平凡。\footnote{回忆我们假设 $G$ 是连通单连通的代数群，此时 $\operatorname{Gr}$ 只有一个连通分支。一般情形下，范畴 $\mathcal{D}_{\kappa_c}\text{-mod}$ 的中心以 $\operatorname{Gr}$ 的连通分支为基，同构于有限群 $\pi_1(G)$ 的群代数。}第二个、也更严重的原因，是范畴 $\mathcal{D}_{\kappa_c}\text{-mod}$ 带有额外的对称性，即 ${}^LG$ 的表示张量范畴 $\operatorname{Rep}{}^LG$ 的作用；而如下所述，此作用在函子 $\Gamma$ 下平凡化。

% [Tutor Note: 此处“中心平凡”依赖 $G$ 单连通（见脚注 6）；若 $G$ 不单连通，$\pi_1(G)$ 会在中心中留下分量。这是把几何 Langlands 对应与经典 Langlands 对应比较时的常见混淆点，务必与 $G$ 的假定一起记忆。]

在 $\operatorname{Op}_{{}^LG}(D)$ 上存在典范的主 ${}^LG$-丛，记作 $\mathcal{P}$。按定义，$\mathcal{P}$ 在 $\chi=(\mathcal{F},\nabla,\mathcal{F}_{{}^LB})\in\operatorname{Op}_{{}^LG}(D)$ 处的纤维是 $\mathcal{F}_0$，即 $\chi$ 底层的 ${}^LG$-丛 $\mathcal{F}$ 在 $0\in D$ 处的纤维。对对象 $V\in\operatorname{Rep}{}^LG$，记 $\mathcal{V}$ 为 $\operatorname{Op}_{{}^LG}(D)$ 上相应的向量丛，即
\[
    \mathcal{V}=\mathcal{P}\times_{{}^LG}V.
\]
其次，考虑 $\operatorname{Gr}$ 上 $G[[t]]$-等变 $\mathcal{D}_{\kappa_c}$-模的范畴 $\mathcal{D}_{\kappa_c}\text{-mod}^{G[[t]]}$。它等价于上面考虑的范畴
\[
    \mathcal{D}\text{-mod}^{G[[t]]}=\mathcal{H}(G((t)),G[[t]]).
\]
这是张量范畴（关于卷积函子），等价于范畴 $\operatorname{Rep}{}^LG$。对 $V\in\operatorname{Rep}{}^LG$ 对应的 $\mathcal{D}_{\kappa_c}\text{-mod}^{G[[t]]}$ 的对象，我们沿用记号 $\mathcal{H}_V$。范畴 $\mathcal{D}_{\kappa_c}\text{-mod}^{G[[t]]}$ 通过卷积函子
\[
    \mathcal{M}\mapsto\mathcal{H}_V\star\mathcal{M}
\]
作用在 $\mathcal{D}_{\kappa_c}\text{-mod}$ 上，这些函子正合。这相当于范畴 $\operatorname{Rep}{}^LG$ 在 $\mathcal{D}_{\kappa_c}\text{-mod}$ 上的张量作用。

现在，由 [Beilinson and Drinfeld 1997] 的结果，存在函子性同构
\[
    \Gamma(\operatorname{Gr},\mathcal{H}_V\star\mathcal{M})
    \simeq\Gamma(\operatorname{Gr},\mathcal{M})\otimes_{\mathfrak{z}}\mathcal{V},
    \qquad V\in\operatorname{Rep}{}^LG,
\]
与张量结构相容（细节见 [Frenkel and Gaitsgory 2006c; Frenkel and Gaitsgory 2005]）。于是我们看到，$\mathcal{D}_{\kappa_c}\text{-mod}$ 中存在不同构的对象，而函子 $\Gamma$ 把它们送到 $\hat{\mathfrak{g}}_{\kappa_c}\text{-mod}_{reg}$ 中同构的对象。因此必须修改范畴 $\mathcal{D}_{\kappa_c}\text{-mod}$ 与函子 $\Gamma$，才有可能得到等价于 $\hat{\mathfrak{g}}_{\kappa_c}\text{-mod}_{reg}$ 的范畴。

在 [Frenkel and Gaitsgory 2006c] 中证明了如何修改范畴 $\mathcal{D}_{\kappa_c}\text{-mod}$：同时向它“添加” $\mathfrak{z}$ 作为中心，并“除以”上述 $\operatorname{Rep}{}^LG$-作用。结果得到一个可能等价于 $\hat{\mathfrak{g}}_{\kappa_c}\text{-mod}_{reg}$ 的范畴的候选者。这就是 $\operatorname{Gr}$ 上的 Hecke 本征模范畴，记作 $\mathcal{D}^{\text{Hecke}}_{\kappa_c}\text{-mod}_{reg}$。

按定义，$\mathcal{D}^{\text{Hecke}}_{\kappa_c}\text{-mod}_{reg}$ 的对象是 $\mathcal{D}_{\kappa_c}\text{-mod}$ 的对象，连同代数 $\mathfrak{z}$ 通过自同态的作用以及一族同构
\[
    \alpha_V:\mathcal{H}_V\star\mathcal{M}\xrightarrow{\ \sim\ }
    \underline{V}\otimes_{\mathfrak{z}}\mathcal{M},
    \qquad V\in\operatorname{Rep}{}^LG,
\]
且与张量结构相容。

上述函子 $\Gamma$ 自然地给出函子
\begin{equation}
    \Gamma^{\text{Hecke}}:\mathcal{D}^{\text{Hecke}}_{\kappa_c}\text{-mod}_{reg}
    \to\hat{\mathfrak{g}}_{\kappa_c}\text{-mod}_{reg}.
    \tag{10.3-22}
\end{equation}

这其实是一般性质。为说明这点，设我们有一个阿贝尔范畴 $\mathcal{C}$，张量范畴 $\operatorname{Rep}H$ 作用其上，其中 $H$ 是代数群；记此作用为
\[
    \mathcal{M}\mapsto\mathcal{M}\star V,\qquad V\in\operatorname{Rep}H.
\]
设 $\mathcal{C}^{\text{Hecke}}$ 是这样的范畴：其对象是集合 $(\mathcal{M},\{\alpha_V\}_{V\in\operatorname{Rep}H})$，其中 $\mathcal{M}\in\mathcal{C}$，而 $\{\alpha_V\}$ 是一族相容的同构
\[
    \alpha_V:\mathcal{M}\star V\xrightarrow{\ \sim\ }
    \underline{V}\otimes_{\mathbb{C}}\mathcal{M},
    \qquad V\in\operatorname{Rep}H,
\]
其中 $\underline{V}$ 是 $V$ 底层的向量空间。范畴 $\mathcal{C}^{\text{Hecke}}$ 带有群 $H$ 的自然作用：对 $h\in H$，令
\[
    h\cdot(\mathcal{M},\{\alpha_V\}_{V\in\operatorname{Rep}H})
    =(\mathcal{M},\{(h\otimes\operatorname{id}_{\mathcal{M}})\circ\alpha_V\}_{V\in\operatorname{Rep}H}).
\]
换言之，$\mathcal{M}$ 保持不变，但同构 $\alpha_V$ 与 $h$ 复合。

范畴 $\mathcal{C}$ 可以重新实现为 $\mathcal{C}^{\text{Hecke}}$ 关于此作用的 $H$-等变对象范畴，见 [Arkhipov and Gaitsgory 2003]。因此可以把 $\mathcal{C}^{\text{Hecke}}$ 看作范畴 $\mathcal{C}$ 的“去等变化”。\transnote{“去等变化”（de-equivariantization）与“等变化”（equivariantization）互为逆操作：给定 $\operatorname{Rep}H$ 在 $\mathcal{C}$ 上的张量作用，$\mathcal{C}^{\text{Hecke}}$ 的对象是带相容同构 $\alpha_V$ 的 $\mathcal{M}$，而 $H$ 通过复合 $\alpha_V$ 作用其上；再取 $H$-等变对象便回到 $\mathcal{C}$。在本书情形 $H={}^LG$、$\mathcal{C}=\mathcal{D}_{\kappa_c}\text{-mod}$，这一步同时把 $\mathfrak{z}$-中心“添加”进来。名字里的“等变化”指取 $H$-等变对象，不是取 $G[[t]]$-等变模，二者不可混淆。[来源：原书 p.331；S. Arkhipov and D. Gaitsgory, Another realization of the category of modules over the small quantum group, Adv. Math. 173 (2003), 114--143.]}

设函子 $G:\mathcal{C}\to\mathcal{C}'$ 满足函子性同构
\begin{equation}
    G(\mathcal{M}\star V)\simeq G(\mathcal{M})\otimes_{\mathbb{C}}\underline{V},
    \qquad V\in\operatorname{Rep}H,
    \tag{10.3-23}
\end{equation}
且与张量结构相容。则据 [Arkhipov and Gaitsgory 2003]，存在函子
\[
    G^{\text{Hecke}}:\mathcal{C}^{\text{Hecke}}\to\mathcal{C}',
\]
使得 $G\simeq G^{\text{Hecke}}\circ\operatorname{Ind}$，其中函子 $\operatorname{Ind}:\mathcal{C}\to\mathcal{C}^{\text{Hecke}}$ 把 $\mathcal{M}$ 送到 $\mathcal{M}\star\mathcal{O}_H$，这里 $\mathcal{O}_H$ 是 $H$ 的正则表示。函子 $G^{\text{Hecke}}$ 可显式描述如下：同构 $\alpha_V$ 与 (10.3-23) 给出代数 $\mathcal{O}_H$ 在 $G(\mathcal{M})$ 上的作用，而 $G^{\text{Hecke}}(\mathcal{M})$ 是取 $G(\mathcal{M})$ 在 $1\in H$ 处的纤维。

在我们的情形，取 $\mathcal{C}=\mathcal{D}_{\kappa_c}\text{-mod}$，$\mathcal{C}'=\hat{\mathfrak{g}}_{\kappa_c}\text{-mod}_{reg}$，$G=\Gamma$。唯一的区别是现在我们在底 $\operatorname{Op}_{{}^LG}(D)$ 上工作，必须把它考虑进去。于是函子 $\Gamma$ 分解为
\[
    \Gamma\simeq\Gamma^{\text{Hecke}}\circ\operatorname{Ind},
\]
其中 $\Gamma^{\text{Hecke}}$ 是函子 (10.3-22)（细节见 [Frenkel and Gaitsgory 2006c; Frenkel and Gaitsgory 2005]）。此外，群 $G((t))$ 在 $\operatorname{Gr}$ 上的左作用给出它在范畴 $\mathcal{D}^{\text{Hecke}}_{\kappa_c}\text{-mod}_{reg}$ 上的作用，而函子 $\Gamma^{\text{Hecke}}$ 把这个作用与 $G((t))$ 在 $\hat{\mathfrak{g}}_{\kappa_c}\text{-mod}_{reg}$ 上的作用交织起来。

% [Tutor Note: “去等变化”（de-equivariantization）是本小节的技术核心：先人为添加 $\mathfrak{z}$-中心作用（从而放大范畴），再除掉 $\operatorname{Rep}{}^LG$ 的对称性。它是 Arkhipov–Gaitsgory 的通用构造在 $\operatorname{Op}_{{}^LG}(D)$ 上的相对版本，理解这一点才能明白为何 $\Gamma$ 可以通过 $\Gamma^{\text{Hecke}}$ 分解。]

下述猜想见 [Frenkel and Gaitsgory 2006c]：

\begin{conjecture}
    公式 (10.3-22) 中的函子 $\Gamma^{\text{Hecke}}$ 给出范畴 $\mathcal{D}^{\text{Hecke}}_{\kappa_c}\text{-mod}_{reg}$ 与 $\hat{\mathfrak{g}}_{\kappa_c}\text{-mod}_{reg}$ 的等价。
\end{conjecture}

[Frenkel and Gaitsgory 2006c] 证明了函子 $\Gamma^{\text{Hecke}}$ 延拓到导出范畴后是全忠实的。此外，[Frenkel and Gaitsgory 2005] 证明了它给出相应的 $I^0$-等变范畴的等价，其中 $I^0=[I,I]$ 是 Iwahori 子群的根。\transnote{Iwahori 子群 $I$ 是 $G[[t]]$ 在标准 Borel 子群 $B\subset G$ 下的逆像；$I^0=[I,I]$ 是它的幂幺根（pro-unipotent radical）。以 $I$ 代替 $G[[t]]$ 得到的 Iwahori Hecke 范畴由仿射旗簇 $\operatorname{Fl}=G((t))/I$ 上 $I$-等变的 $D$-模构成；其卷积并非逐项正合，通常须在导出范畴或混合范畴中处理，这与球面情形的正合性形成对照。正文所引 [Frenkel and Gaitsgory 2005] 的等价是在 $I^0$-等变层面建立的；球面情形保留 $G[[t]]$-等变对象（较粗），Iwahori 情形保留 $I$-等变对象（较细），后者信息更多但代数结构更复杂。[来源：原书 p.332；E. V. Frenkel and D. Gaitsgory, Localization of $\hat{\mathfrak{g}}$-modules on the affine Grassmannian, arXiv:math.RT/0512562.]}

让我们把猜想 10.3.9 特殊化到点 $\chi=(\mathcal{F},\nabla,\mathcal{F}_{{}^LB})\in\operatorname{Op}_{{}^LG}(D)$。则右端考虑范畴 $\hat{\mathfrak{g}}_{\kappa_c}\text{-mod}_\chi$，左端考虑范畴 $\mathcal{D}^{\text{Hecke}}_{\kappa_c}\text{-mod}_\chi$。它的对象由一个 $\mathcal{D}_{\kappa_c}$-模 $\mathcal{M}$ 与一族同构
\begin{equation}
    \alpha_V:\mathcal{H}_V\star\mathcal{M}\xrightarrow{\ \sim\ }
    \underline{V}_{\mathcal{F}_0}\otimes\mathcal{M},
    \qquad V\in\operatorname{Rep}{}^LG
    \tag{10.3-24}
\end{equation}
组成，其中 $\underline{V}_{\mathcal{F}_0}$ 是表示 $V$ 经 ${}^LG$-torsor $\mathcal{F}_0$ 的扭变。这些同构必须与范畴 $\mathcal{H}(G((t)),G[[t]])$ 上的张量结构相容。

猜想 10.3.9 蕴含存在典范的范畴等价
\begin{equation}
    \mathcal{D}^{\text{Hecke}}_{\kappa_c}\text{-mod}_\chi
    \simeq\hat{\mathfrak{g}}_{\kappa_c}\text{-mod}_\chi.
    \tag{10.3-25}
\end{equation}
正是这个猜想中的等价应被看作 Beilinson--Bernstein 等价的类比。

从此观点看，每个范畴 $\mathcal{D}^{\text{Hecke}}_{\kappa_c}\text{-mod}_\chi$ 都可以看作对应于平凡 ${}^LG$-局部系统的所求 Langlands 范畴 $\mathcal{C}_{\sigma_0}$ 的第二个化身。

现在我们对“为何宜把 $\mathcal{D}^{\text{Hecke}}_{\kappa_c}\text{-mod}_\chi$ 看作群 $G(F)$ 无分歧表示的范畴化”给出另一种解释。首先注意，这些范畴彼此等价，且都等价于范畴 $\mathcal{D}^{\text{Hecke}}_{\kappa_c}\text{-mod}$，后者的对象是 $\mathcal{D}_{\kappa_c}$-模 $\mathcal{M}$ 连同一族同构
\begin{equation}
    \alpha_V:\mathcal{H}_V\star\mathcal{M}\xrightarrow{\ \sim\ }
    \underline{V}\otimes\mathcal{M},
    \qquad V\in\operatorname{Rep}{}^LG.
    \tag{10.3-26}
\end{equation}
比较公式 (10.3-24) 与 (10.3-26) 可见，对 ${}^LG$-torsor $\mathcal{F}_0$ 的每个平凡化选择\transnote{不同平凡化相差 $\mathcal{F}_0$ 的一个自同构，即 ${}^LG$ 的一个元素；把这个元素作用在 $\alpha_V$ 上，就得到另一个等价。因此 (10.3-24) 与 (10.3-26) 的范畴只在选定平凡化后才等同，而所有选择的全体构成 $\mathcal{F}_0$ 上的一个 torsor。这正是下一小节群胚 $\operatorname{Isom}$ 作用的来源：不是范畴本身典范等价，而是“带扭变数据的范畴”之间的等价。[来源：原书 p.333；E. Frenkel and D. Gaitsgory, Local geometric Langlands correspondence and affine Kac-Moody algebras, 2006, §2.]}（$\mathcal{F}_0$ 是 oper $\chi$ 底层的 $D$ 上主 ${}^LG$-丛 $\mathcal{F}$ 在 $0\in D$ 处的纤维），有等价
\[
    \mathcal{D}^{\text{Hecke}}_{\kappa_c}\text{-mod}_\chi
    \simeq\mathcal{D}^{\text{Hecke}}_{\kappa_c}\text{-mod}.
\]

% [Tutor Note: (10.3-24) 与 (10.3-26) 的唯一差别是是否用 ${}^LG$-torsor $\mathcal{F}_0$ 扭变 $V$。因此两者只在“选定了 $\mathcal{F}_0$ 的平凡化”后才等同；不同平凡化给出不同的等价，这正是 10.3.7 节群胚 $\operatorname{Isom}$ 作用的来源，也是“几何情形中 ${}^LG$ 的作用取代共轭类 $\gamma$”的具体机制。]

现在回忆 10.3.1 节：由 Satake 对应 (10.3-2)，${}^LG$ 中每个半单共轭类 $\gamma$ 对应于 $G(F)$ 的不可约无分歧表示 $(R_\gamma,\pi_\gamma)$。已知存在非退化配对
\[
    \langle,\rangle:R_\gamma\times R_{\gamma^{-1}}\to\mathbb{C},
\]
换言之，$R_{\gamma^{-1}}$ 是 $G(F)$ 的与 $R_\gamma$ 反变的表示（它可以实现在 $R_\gamma$ 的对偶空间的光滑向量空间中）。

设 $v\in R_{\gamma^{-1}}$ 是非零向量且 $K_0v=v$（此向量在相差一个标量意义下唯一）。则它满足 Hecke 本征向量性质 (10.3-5)（其中需把 $\gamma$ 换成 $\gamma^{-1}$）。这允许我们利用矩阵系数把 $R_\gamma$ 嵌入 $G(F)$ 上光滑局部常值右 $K_0$-不变函数空间（等价地，$G(F)/K_0$ 上的函数空间）：
\[
    u\in R_\gamma\mapsto f_u,\qquad f_u(g)=\langle u,gv\rangle.
\]
Hecke 本征向量性质 (10.3-5) 蕴含函数 $f_u$ 右 $K_0$-不变，且满足
\begin{equation}
    f\star H_V=\operatorname{Tr}(\gamma^{-1},V)f,
    \tag{10.3-27}
\end{equation}
其中 $\star$ 表示卷积积 (10.2-1)。设 $C(G(F)/K_0)_\gamma$ 是 $G(F)/K_0$ 上满足 (10.3-27) 的光滑局部常值函数空间。它带有由左作用在 $G(F)/K_0$ 上诱导的 $G(F)$ 表示。我们构造了单射 $R_\gamma\to C(G(F)/K_0)_\gamma$，且可以证明对一般的 $\gamma$ 它是同构。

于是我们得到 $G(F)$ 的不可约无分歧表示在满足 Hecke 本征函数条件 (10.3-27) 的商 $G(F)/K_0$ 上的函数空间中的实现。Hecke 本征模条件 (10.3-26) 可看作 (10.3-27) 的范畴类比。因此，$\operatorname{Gr}=G((t))/K_0$ 上满足 Hecke 本征模条件 (10.3-26) 的扭变 $D$-模范畴 $\mathcal{D}^{\text{Hecke}}_{\kappa_c}\text{-mod}$，连同其 $G((t))$-作用，似乎是 $G(F)$ 的不可约无分歧表示的自然的范畴化。

\subsection{模范畴之间的等价}

$\operatorname{Op}_{{}^LG}(D)$ 中所有 opers 都对应同一个 ${}^LG$-局部系统，即平凡局部系统。因此按猜想 10.1.2，我们预期范畴 $\hat{\mathfrak{g}}_{\kappa_c}\text{-mod}_\chi$ 彼此等价。更确切地说，对 $\operatorname{Op}_{{}^LG}(D)$ 中任意两个 opers 的底层局部系统之间的每个同构，我们希望有相应范畴之间的等价，且这些等价应与同构的复合运算相容。

让我们详细说明。设 $\chi=(\mathcal{F},\nabla,\mathcal{F}_{{}^LB})$ 与 $\chi'=(\mathcal{F}',\nabla',\mathcal{F}'_{{}^LB})$ 是 $\operatorname{Op}_{{}^LG}(D)$ 中的两个 opers。则底层局部系统之间的同构 $(\mathcal{F},\nabla)\xrightarrow{\sim}(\mathcal{F}',\nabla')$ 等同于 ${}^LG$-torsors $\mathcal{F}_0$ 与 $\mathcal{F}'_0$ 之间的同构 $\mathcal{F}_0\xrightarrow{\sim}\mathcal{F}'_0$，其中 $\mathcal{F}_0$、$\mathcal{F}'_0$ 分别是 ${}^LG$-丛 $\mathcal{F}$、$\mathcal{F}'$ 在 $0\in D$ 处的纤维。记这族同构为 $\operatorname{Isom}_{\chi,\chi'}$。则有
\[
    \operatorname{Isom}_{\chi,\chi'}
    =\mathcal{F}_0\times_{{}^LG}{}^LG\times_{{}^LG}\mathcal{F}'_0,
\]
其中我们分别关于左作用以 $\mathcal{F}_0$、关于右作用以 $\mathcal{F}'_0$ 扭变 ${}^LG$。特别地，
\[
    \operatorname{Isom}_{\chi,\chi}
    ={}^LG_{\mathcal{F}_0}=\mathcal{F}_0\times_{{}^LG}\operatorname{Ad}{}^LG
\]
正是 $\mathcal{F}_0$ 的自动同构群。

把集合 $\operatorname{Isom}_{\chi,\chi'}$ 合并为 $\operatorname{Op}_{{}^LG}(D)$ 上的群胚 $\operatorname{Isom}$ 是有启发性的。按定义，$\operatorname{Isom}$ 由三元组 $(\chi,\chi',\phi)$ 组成，其中 $\chi,\chi'\in\operatorname{Op}_{{}^LG}(D)$，$\phi\in\operatorname{Isom}_{\chi,\chi'}$ 是底层局部系统的同构。两个态射 $\operatorname{Isom}\to\operatorname{Op}_{{}^LG}(D)$ 分别把三元组送到 $\chi$ 与 $\chi'$。恒等态射 $\operatorname{Op}_{{}^LG}(D)\to\operatorname{Isom}$ 把 $\chi$ 送到 $(\chi,\chi,\operatorname{Id})$，而复合态射
\[
    \operatorname{Isom}\times_{\operatorname{Op}_{{}^LG}(D)}\operatorname{Isom}
    \to\operatorname{Isom}
\]
对应于复合两个同构。

猜想 10.1.2 对正则 opers 有如下更精确的表述。

\begin{conjecture}
    对每个 $\phi\in\operatorname{Isom}_{\chi,\chi'}$，存在等价
    \[
        E_\phi:\hat{\mathfrak{g}}_{\kappa_c}\text{-mod}_\chi
        \to\hat{\mathfrak{g}}_{\kappa_c}\text{-mod}_{\chi'},
    \]
    它交织两个范畴上 $G((t))$ 的作用，使得 $E_{\operatorname{Id}}=\operatorname{Id}$，且存在同构
    $\beta_{\phi,\phi'}:E_{\phi\circ\phi'}\simeq E_\phi\circ E_{\phi'}$，满足
    \[
        \beta_{\phi\circ\phi',\phi''}\beta_{\phi,\phi'}
        =\beta_{\phi,\phi'\circ\phi''}\beta_{\phi',\phi''}
    \]
    对所有可在适当顺序下复合的同构 $\phi,\phi',\phi''$ 成立。
\end{conjecture}

% [Tutor Note: 这一相容性条件就是群胚作用的结合律：$\beta$ 有两组指标（两个可复合的态射），而四个 $\beta$ 因子的排列顺序极易抄错。核对时可按“先复合哪一对”来记：左式为 $(\phi\circ\phi')\circ\phi''$ 一侧，右式为 $\phi\circ(\phi'\circ\phi'')$ 一侧。]

也就是说，$\operatorname{Op}_{{}^LG}(D)$ 上的群胚 $\operatorname{Isom}$ 作用在以 $\operatorname{Op}_{{}^LG}(D)$ 为底、纤维范畴为 $\hat{\mathfrak{g}}_{\kappa_c}\text{-mod}_\chi$ 的范畴 $\hat{\mathfrak{g}}_{\kappa_c}\text{-mod}_{reg}$ 上，并沿纤维保持 $G((t))$ 的作用。

特别地，此猜想蕴含群 ${}^LG_{\mathcal{F}_0}$ 作用在范畴 $\hat{\mathfrak{g}}_{\kappa_c}\text{-mod}_\chi$ 上，对任意 $\chi\in\operatorname{Op}_{{}^LG}(D)$。

现在注意，猜想 10.3.9 蕴含猜想 10.3.10。事实上，由猜想 10.3.9，有典范的范畴等价 (10.3-25)。由范畴 $\mathcal{D}^{\text{Hecke}}_{\kappa_c}\text{-mod}_\chi$ 的定义（即 (10.3-24)）可知，对每个同构 $\phi\in\operatorname{Isom}_{\chi,\chi'}$——即 opers $\chi$ 与 $\chi'$ 底层的 ${}^LG$-torsors $\mathcal{F}_0$ 与 $\mathcal{F}'_0$ 的同构——有典范等价
\[
    \mathcal{D}^{\text{Hecke}}_{\kappa_c}\text{-mod}_\chi
    \simeq\mathcal{D}^{\text{Hecke}}_{\kappa_c}\text{-mod}_{\chi'}.
\]
因此得到所求的等价
\[
    E_\phi:\hat{\mathfrak{g}}_{\kappa_c}\text{-mod}_\chi
    \to\hat{\mathfrak{g}}_{\kappa_c}\text{-mod}_{\chi'}.
\]
此外，显然这些等价满足猜想 10.3.10 的条件。特别地，它们交织 $G((t))$ 的作用，因为 $G((t))$ 的作用影响 $\mathcal{D}^{\text{Hecke}}_{\kappa_c}\text{-mod}_\chi$ 的对象底层的 $D$-模 $\mathcal{M}$，但不影响同构 $\alpha_V$。

等价地，我们可以说群胚 $\operatorname{Isom}$ 自然地作用在范畴 $\mathcal{D}^{\text{Hecke}}_{\kappa_c}\text{-mod}_{reg}$ 上。由猜想 10.3.9，这给出 $\operatorname{Isom}$ 在 $\hat{\mathfrak{g}}_{\kappa_c}\text{-mod}_{reg}$ 上的作用。

特别地，我们构造了群 $({}^LG)_{\mathcal{F}_0}$（${}^LG$ 经某个 oper $\chi$ 底层的 ${}^LG$-torsor $\mathcal{F}_0$ 的扭变）在范畴 $\mathcal{D}^{\text{Hecke}}_{\kappa_c}\text{-mod}_\chi$ 上的作用。事实上，每个元素 $g\in({}^LG)_{\mathcal{F}_0}$ 作用在 ${}^LG$ 的任何有限维表示 $V$ 的 $\mathcal{F}_0$-扭变 $V_{\mathcal{F}_0}$ 上。给定 $\mathcal{D}^{\text{Hecke}}_{\kappa_c}\text{-mod}_{\chi'}$ 的对象 $(\mathcal{M},(\alpha_V))$\transnote{此处原文印作 $\chi'$，但本段正在构造群 $({}^LG)_{\mathcal{F}_0}$ 对范畴 $\mathcal{D}^{\text{Hecke}}_{\kappa_c}\text{-mod}_\chi$ 的作用：新对象仍在同一 $\chi$ 处（只把 $\alpha_V$ 与 $g$ 复合），若读作 $\chi'$ 则作用跑到另一个范畴去了。故 $\chi'$ 应为 $\chi$ 的笔误；译文照录原文，读者按 $\chi$ 理解即可。[来源：原书 p.335；对照同页上一段“on the category $\mathcal{D}^{\text{Hecke}}_{\kappa_c}\text{-mod}_\chi$”。]}，我们构造新对象 $(\mathcal{M},((g\otimes\operatorname{Id}_{\mathcal{M}})\circ\alpha_V))$。这样我们不改变 $D$-模 $\mathcal{M}$，但通过把 Hecke 本征模条件 (10.3-24) 中出现的同构 $\alpha_V$ 与 $g$ 在 $V_{\mathcal{F}_0}$ 上的作用复合来改变它们。由猜想 10.3.9，范畴 $\mathcal{D}^{\text{Hecke}}_{\kappa_c}\text{-mod}_\chi$ 等价于 $\hat{\mathfrak{g}}_{\kappa_c}\text{-mod}_\chi$。因此这给出群 $({}^LG)_{\mathcal{F}_0}$ 在 $\hat{\mathfrak{g}}_{\kappa_c}\text{-mod}_\chi$ 上的作用。但用 $\hat{\mathfrak{g}}_{\kappa_c}$-模来表述这个作用要困难得多。下文将看到这些对称作用的例子。

\subsection{推广到其他整优权}

上面我们详尽研究了对应于正则 opers $\chi\in\operatorname{Op}_{{}^LG}(D)$ 的范畴 $\hat{\mathfrak{g}}_{\kappa_c}\text{-mod}_\chi$ 与 $\hat{\mathfrak{g}}_{\kappa_c}\text{-mod}^{G[[t]]}_\chi$。然而由引理 10.3.1，映射 $\alpha:\operatorname{Op}_{{}^LG}(D^\times)\to\operatorname{Loc}_{{}^LG}(D^\times)$ 在平凡局部系统 $\sigma_0$ 上的（集合论）纤维，是子集 $\operatorname{Op}^\lambda_{{}^LG}$（$\lambda\in P^+$）的不交并。这里我们简要讨论 $\chi\in\operatorname{Op}^\lambda_{{}^LG}$（$\lambda\neq0$）时的范畴 $\hat{\mathfrak{g}}_{\kappa_c}\text{-mod}_\chi$ 与 $\hat{\mathfrak{g}}_{\kappa_c}\text{-mod}^{G[[t]]}_\chi$。

考虑 9.6 节定义的以 $\lambda$ 为最高权的 Weyl 模 $\mathbb{V}_\lambda$。由定理 9.6.1，有
\begin{equation}
    \operatorname{End}_{\hat{\mathfrak{g}}_{\kappa_c}}\mathbb{V}_\lambda
    \simeq\operatorname{Fun}\operatorname{Op}^\lambda_{{}^LG}.
    \tag{10.3-28}
\end{equation}
设 $\chi\in\operatorname{Op}^\lambda_{{}^LG}\subset\operatorname{Op}_{{}^LG}(D^\times)$。则 $\chi$ 定义了代数 $\operatorname{End}_{\hat{\mathfrak{g}}_{\kappa_c}}\mathbb{V}_\lambda$ 的特征标。设 $\mathbb{V}_\lambda(\chi)$ 是 $\mathbb{V}_\lambda$ 对该特征标之核的商。下述 [Frenkel and Gaitsgory 2007c] 的猜想\footnote{校样时补注：猜想 10.3.11–10.3.13 已由 E. Frenkel 与 D. Gaitsgory 在 [Frenkel and Gaitsgory 2007b] 中证明。}是定理 10.3.3 的类比：

\begin{conjecture}
    设 $\chi\in\operatorname{Op}^\lambda_{{}^LG}\subset\operatorname{Op}_{{}^LG}(D^\times)$。则范畴 $\hat{\mathfrak{g}}_{\kappa_c}\text{-mod}^{G[[t]]}_\chi$ 等价于向量空间范畴：它在同构意义下唯一的不可约对象是 $\mathbb{V}_\lambda(\chi)$，任何其他对象都同构于 $\mathbb{V}_\lambda(\chi)$ 的若干拷贝的直和。\transnote{本节的最后三条猜想构成一组：10.3.11 断言对任意 $\chi\in\operatorname{Op}^\lambda$（$\lambda\in P^+$），$G[[t]]$-等变范畴仍等价于向量空间范畴，且唯一不可约对象是 Weyl 模 $\mathbb{V}_\lambda(\chi)$；10.3.12 进一步断言 $F_\lambda$ 与 $G_\lambda$ 在整个族 $\operatorname{Op}^\lambda$ 上互逆；10.3.13 则给出两条判据——$\mathbb{V}_\lambda$ 在 $\mathfrak{z}_\lambda$ 上自由，以及 $\operatorname{Ext}^i\simeq\Omega^i_{\mathfrak{z}_\lambda}$（(10.3-29)）。脚注 7 已说明：三条均由 Frenkel--Gaitsgory 在 [Frenkel and Gaitsgory 2007b] 中证明，该文标题正是“球面情形”。[来源：原书 p.336--337；E. Frenkel and D. Gaitsgory, Local geometric Langlands correspondence: the spherical case, arXiv:0711.1132.]}
\end{conjecture}

注意这与猜想 10.1.2 一致：后者告诉我们，对一切与 $D$ 上平凡局部系统规范等价的 opers，范畴 $\hat{\mathfrak{g}}_{\kappa_c}\text{-mod}^{G[[t]]}_\chi$ 应彼此等价。

与 $\lambda=0$ 的情形一样，考虑整个族 $\operatorname{Op}^\lambda_{{}^LG}$ 而非单个 ${}^LG$-oper $\chi$ 更方便。设 $\hat{\mathfrak{g}}_{\kappa_c}\text{-mod}_{\lambda,reg}$ 是范畴 $\hat{\mathfrak{g}}_{\kappa_c}\text{-mod}$ 的满子范畴，其对象是那些中心 $Z(\hat{\mathfrak{g}})$ 的作用通过同态
\[
    Z(\hat{\mathfrak{g}})\simeq\operatorname{Fun}\operatorname{Op}_{{}^LG}(D^\times)
    \to\mathfrak{z}_\lambda
\]
分解的 $\hat{\mathfrak{g}}_{\kappa_c}$-模，其中我们令
\[
    \mathfrak{z}_\lambda=\operatorname{Fun}\operatorname{Op}^\lambda_{{}^LG}.
\]
设 $\hat{\mathfrak{g}}_{\kappa_c}\text{-mod}^{G[[t]]}_{\lambda,reg}$ 是相应的 $G[[t]]$-等变范畴。于是
\[
    \hat{\mathfrak{g}}_{\kappa_c}\text{-mod}_{\lambda,reg}
    \quad\text{与}\quad
    \hat{\mathfrak{g}}_{\kappa_c}\text{-mod}^{G[[t]]}_{\lambda,reg}
\]
是以 $\operatorname{Op}^\lambda_{{}^LG}$ 为底的纤维范畴，在 $\chi\in\operatorname{Op}^\lambda_{{}^LG}$ 上的纤维分别是 $\hat{\mathfrak{g}}_{\kappa_c}\text{-mod}_\chi$ 与 $\hat{\mathfrak{g}}_{\kappa_c}\text{-mod}^{G[[t]]}_\chi$。

关于范畴 $\hat{\mathfrak{g}}_{\kappa_c}\text{-mod}^{G[[t]]}_{\lambda,reg}$，我们有下述猜想式的描述；它蕴含猜想 10.3.11（见 [Frenkel and Gaitsgory 2007b]）。

设 $\mathfrak{z}_\lambda\text{-mod}$ 是交换代数 $\mathfrak{z}_\lambda$ 上的模范畴。等价地，它是空间 $\operatorname{Op}^\lambda_{{}^LG}$ 上的拟凝聚层范畴。

按定义，$\hat{\mathfrak{g}}_{\kappa_c}\text{-mod}^{G[[t]]}_{\lambda,reg}$ 的任何对象都是 $\mathfrak{z}_\lambda$-模。引入函子
\[
    F_\lambda:\hat{\mathfrak{g}}_{\kappa_c}\text{-mod}^{G[[t]]}_{\lambda,reg}
    \to\mathfrak{z}_\lambda\text{-mod},
    \qquad M\mapsto\operatorname{Hom}_{\hat{\mathfrak{g}}_{\kappa_c}}(\mathbb{V}_\lambda,M),
\]
\[
    G_\lambda:\mathfrak{z}_\lambda\text{-mod}
    \to\hat{\mathfrak{g}}_{\kappa_c}\text{-mod}^{G[[t]]}_{\lambda,reg},
    \qquad \mathcal{F}\mapsto\mathbb{V}_\lambda\otimes_{\mathfrak{z}_\lambda}\mathcal{F}.
\]

\begin{conjecture}
    函子 $F_\lambda$ 与 $G_\lambda$ 是互逆的范畴等价。
\end{conjecture}

我们预期此猜想可以沿 10.3.4 节给出的定理 10.3.4 的证明思路来证明。更确切地说，它应当由 [Frenkel and Gaitsgory 2007c] 的猜想推出，后者推广了 $\lambda=0$ 时的相应陈述。

\begin{conjecture}
    （1）$\mathbb{V}_\lambda$ 作为 $\mathfrak{z}_\lambda$-模是自由的。
    （2）
    \begin{equation}
        \operatorname{Ext}^i_{\hat{\mathfrak{g}},G[[t]]}(\mathbb{V}_\lambda,\mathbb{V}_\lambda)
        \simeq\Omega^i_{\mathfrak{z}_\lambda},
        \tag{10.3-29}
    \end{equation}
    即 $\operatorname{Op}^\lambda_{{}^LG}$ 上 $i$ 次微分形式构成的空间。
\end{conjecture}

% [Tutor Note: 原书 p.337 把猜想 10.3.13 的排版做成 “猜想 10.3.13. (1) …  (2)” 后接带编号 (10.3-29) 的显示公式，再以斜体文字收尾；译文按语义把 (2) 与公式放在同一项内。注意 (10.3-29) 中上同调的下标是 $\hat{\mathfrak{g}}$（未固定水平），与定理 10.3.7 的约定一致，而不是 $\hat{\mathfrak{g}}_{\kappa_c}$。]

% 本节（§10.3）至此结束。原书 §10.4「驯顺分歧情形」自 PDF p.347（印刷 p.337）下半页开始，属 U16（ch10-3.tex）。



